\clearpage
\section*{Problem 4}
\addcontentsline{toc}{section}{Problem 4}
\begingroup
\hypersetup{colorlinks=true,
            linkcolor=blue!50!black,
            citecolor=blue!50!black,
            urlcolor=blue!50!black}
\newcommand{\qfourR}{\mathbb{R}}
\newcommand{\qfourZ}{\mathbb{Z}}
\newcommand{\qfourbn}{\boxplus_n}
\newcommand{\qfourScore}{\mathsf{S}^{\mathrm{ff}}_n}
\newcommand{\qfourInfo}{\mathcal{K}^{\mathrm{ff}}_n}
\newcommand{\qfourRootMap}{\Psi_{\qfourbn}}
\newcommand{\qfourJac}{\mathsf{D}_{\qfourbn}}
\newcommand{\qfourVar}{\operatorname{Var}}
\newcommand{\qfourHess}{\operatorname{Hess}}
\newcommand{\qfourPSD}{\succeq 0}
\newcommand{\qfoursgn}{\operatorname{sgn}}
\newtheorem{qfourtheorem}{Theorem}
\newtheorem{qfourlemma}{Lemma}
\newtheorem{qfourproposition}{Proposition}
\newtheorem{qfourcorollary}{Corollary}
\theoremstyle{definition}
\newtheorem{qfourobservation}{Observation}
\newtheorem{qfourdefinition}{Definition}
\title{Solution to Problem 4: The Finite Free Stam Inequality}
\author{}
\date{}
\maketitle


\section*{Statement and Answer}

For a monic real-rooted polynomial $p$ of degree $n$ with roots
$\alpha=(\alpha_1,\dots,\alpha_n)$, define the \emph{finite free score vector}
$\qfourScore(\alpha)\in(\qfourR\cup\{\infty\})^n$ by
\[
\qfourScore(\alpha)_i \;=\; \sum_{j\ne i}\frac{1}{\alpha_i-\alpha_j},
\]
and the \emph{finite free Fisher information}
$\qfourInfo(p):=\|\qfourScore(\alpha)\|^2$, with $\qfourInfo(p)=\infty$ if $p$ has a
multiple root. Let $\qfourbn$ denote the finite free convolution of degree-$n$
polynomials.

\medskip
\noindent\textbf{Theorem.} \emph{For any two monic real-rooted polynomials
$p,q$ of degree $n$,}
\[
\frac{1}{\qfourInfo(p\qfourbn q)}\;\ge\;\frac{1}{\qfourInfo(p)}+\frac{1}{\qfourInfo(q)}.
\]

\medskip
\noindent The answer is \textbf{YES}. The result was conjectured by
D.~Shlyakhtenko; the proof below follows Blachman's strategy for the classical
Stam inequality, with the key analytic input replaced by a convexity statement
for hyperbolic polynomials.

\section*{Preliminaries on $\qfourbn$}

Recall that $\qfourbn$ is \emph{defined} by the bilinear coefficient formula: if
$p(x)=\sum_{k=0}^n a_kx^{n-k}$ and $q(x)=\sum_{k=0}^n b_kx^{n-k}$ are monic
($a_0=b_0=1$), then $(p\qfourbn q)(x)=\sum_{k=0}^n c_kx^{n-k}$ with
\begin{equation}\label{eq:coeff}
c_k\;=\;\sum_{i+j=k}\frac{(n-i)!\,(n-j)!}{n!\,(n-k)!}\,a_i\,b_j ,
\qquad k=0,1,\dots,n .
\end{equation}
Almost every argument below is more naturally phrased through the equivalent
\emph{symmetrized-product} (Walsh) representation in terms of the roots. We first
establish, with a self-contained combinatorial proof, that the two descriptions
coincide; this is the identity attributed to \cite{MSS22} but which we do not
merely cite.

\begin{qfourlemma}[coefficient formula $=$ symmetrized product]\label{lem:walsh-eq}
Let $p(x)=\prod_{i=1}^n(x-\alpha_i)$ and $q(x)=\prod_{j=1}^n(x-\beta_j)$ be monic
of degree $n$, with roots $\alpha=(\alpha_i)$, $\beta=(\beta_j)$ (real or
complex). Then the polynomial $p\qfourbn q$ defined by \eqref{eq:coeff} satisfies
\begin{equation}\label{eq:walsh}
(p\qfourbn q)(x)\;=\;\frac{1}{n!}\sum_{\pi\in S_n}\prod_{i=1}^n\bigl(x-\alpha_i-\beta_{\pi(i)}\bigr).
\end{equation}
\end{qfourlemma}

\begin{proof}
Throughout, $e_r$ denotes the $r$-th elementary symmetric polynomial, with the
conventions $e_0=1$ and $e_r=0$ for $r<0$ or $r>n$. Since
$p(x)=\prod_i(x-\alpha_i)=\sum_{r}(-1)^r e_r(\alpha)\,x^{n-r}$, the coefficients
of $p$ and $q$ are
\begin{equation}\label{eq:vieta}
a_i=(-1)^i e_i(\alpha),\qquad b_j=(-1)^j e_j(\beta).
\end{equation}

\smallskip\noindent\emph{Step 1: a permutation-averaging identity.}
We claim that for every fixed subset $B\subseteq[n]:=\{1,\dots,n\}$ with $|B|=j$,
\begin{equation}\label{eq:avg-beta}
\frac{1}{n!}\sum_{\pi\in S_n}\;\prod_{b\in B}\beta_{\pi(b)}
\;=\;\frac{j!\,(n-j)!}{n!}\;e_j(\beta),
\end{equation}
in particular the left-hand side does not depend on which set $B$ of size $j$ is
chosen. Indeed, as $\pi$ ranges over $S_n$, its restriction $\pi|_B$ ranges over
all injections $B\hookrightarrow[n]$, and each injection arises from exactly
$(n-j)!$ permutations (the $n-j$ values outside $\pi(B)$ may be assigned to the
$n-j$ indices outside $B$ in $(n-j)!$ ways). Hence
\[
\sum_{\pi\in S_n}\prod_{b\in B}\beta_{\pi(b)}
=(n-j)!\!\!\sum_{\substack{f:B\hookrightarrow[n]\\ \text{injective}}}\;\prod_{b\in B}\beta_{f(b)} .
\]
Grouping the injections $f$ by their image $K=f(B)$ (a $j$-subset of $[n]$), and
noting that for each fixed image there are $j!$ bijections $B\to K$, all giving
the same product $\prod_{k\in K}\beta_k$, the inner sum equals
$j!\sum_{|K|=j}\prod_{k\in K}\beta_k=j!\,e_j(\beta)$. Dividing by $n!$ yields
\eqref{eq:avg-beta}.

\smallskip\noindent\emph{Step 2: expanding the symmetrized product.}
Write $\gamma_i^\pi:=\alpha_i+\beta_{\pi(i)}$. For each fixed $\pi$,
\[
\prod_{i=1}^n\bigl(x-\alpha_i-\beta_{\pi(i)}\bigr)
=\prod_{i=1}^n(x-\gamma_i^\pi)
=\sum_{k=0}^n(-1)^k e_k\bigl(\gamma^\pi\bigr)\,x^{n-k},
\]
so the coefficient of $x^{n-k}$ in the right-hand side of \eqref{eq:walsh} is
\begin{equation}\label{eq:rhs-coeff}
(-1)^k\,\frac1{n!}\sum_{\pi\in S_n} e_k\bigl(\gamma^\pi\bigr).
\end{equation}
Now expand $e_k(\gamma^\pi)=\sum_{|S|=k}\prod_{i\in S}(\alpha_i+\beta_{\pi(i)})$,
where $S$ ranges over $k$-subsets of $[n]$. Multiplying out each factor and
collecting the choices of "$\alpha$" versus "$\beta$" by the index set
$A\subseteq S$ on which $\alpha_i$ is taken (so $\beta_{\pi(i)}$ is taken on
$B:=S\setminus A$), we get
\[
e_k(\gamma^\pi)
=\sum_{|S|=k}\;\sum_{A\subseteq S}\;\Bigl(\prod_{a\in A}\alpha_a\Bigr)\Bigl(\prod_{b\in S\setminus A}\beta_{\pi(b)}\Bigr)
=\sum_{i+j=k}\;\sum_{\substack{A,B\subseteq[n]\ \text{disjoint}\\ |A|=i,\ |B|=j}}
\Bigl(\prod_{a\in A}\alpha_a\Bigr)\Bigl(\prod_{b\in B}\beta_{\pi(b)}\Bigr),
\]
the last equality because choosing a $k$-set $S$ together with a split
$S=A\sqcup B$ is the same as choosing an ordered pair of disjoint sets $(A,B)$
with $|A|=i$, $|B|=j$ and $i+j=k$.

\smallskip\noindent\emph{Step 3: averaging.}
Average over $\pi$ and apply Step~1. The factor $\prod_{a\in A}\alpha_a$ does not
involve $\pi$, and by \eqref{eq:avg-beta} the $\pi$-average of
$\prod_{b\in B}\beta_{\pi(b)}$ equals $\tfrac{j!\,(n-j)!}{n!}e_j(\beta)$,
independently of the particular set $B$ (and of $A$). Therefore, for each split
$i+j=k$,
\[
\frac1{n!}\sum_{\pi\in S_n}\;\sum_{\substack{A,B\ \text{disjoint}\\ |A|=i,\,|B|=j}}
\Bigl(\prod_{a\in A}\alpha_a\Bigr)\Bigl(\prod_{b\in B}\beta_{\pi(b)}\Bigr)
=\frac{j!\,(n-j)!}{n!}\,e_j(\beta)\;
\sum_{\substack{A,B\ \text{disjoint}\\ |A|=i,\,|B|=j}}\prod_{a\in A}\alpha_a .
\]
In the remaining sum the set $B$ is now a free choice of a $j$-subset of the
$n-i$ indices outside $A$, contributing a multiplicity $\binom{n-i}{j}$; hence
\[
\sum_{\substack{A,B\ \text{disjoint}\\ |A|=i,\,|B|=j}}\prod_{a\in A}\alpha_a
=\binom{n-i}{j}\sum_{|A|=i}\prod_{a\in A}\alpha_a
=\binom{n-i}{j}\,e_i(\alpha).
\]
Combining, the $\pi$-average of $e_k(\gamma^\pi)$ is
\[
\frac1{n!}\sum_{\pi}e_k(\gamma^\pi)
=\sum_{i+j=k}\frac{j!\,(n-j)!}{n!}\binom{n-i}{j}\,e_i(\alpha)\,e_j(\beta).
\]
The combinatorial constant simplifies: since
$\binom{n-i}{j}=\dfrac{(n-i)!}{j!\,(n-i-j)!}=\dfrac{(n-i)!}{j!\,(n-k)!}$ (using
$i+j=k$),
\[
\frac{j!\,(n-j)!}{n!}\binom{n-i}{j}
=\frac{j!\,(n-j)!}{n!}\cdot\frac{(n-i)!}{j!\,(n-k)!}
=\frac{(n-i)!\,(n-j)!}{n!\,(n-k)!}.
\]
Hence
\[
\frac1{n!}\sum_{\pi}e_k(\gamma^\pi)
=\sum_{i+j=k}\frac{(n-i)!\,(n-j)!}{n!\,(n-k)!}\,e_i(\alpha)\,e_j(\beta).
\]

\smallskip\noindent\emph{Step 4: matching signs.}
Multiplying by $(-1)^k$ as in \eqref{eq:rhs-coeff} and using $(-1)^k=(-1)^i(-1)^j$
together with \eqref{eq:vieta} (so
$(-1)^i e_i(\alpha)=a_i$ and $(-1)^j e_j(\beta)=b_j$), the coefficient of
$x^{n-k}$ in the right-hand side of \eqref{eq:walsh} equals
\[
\sum_{i+j=k}\frac{(n-i)!\,(n-j)!}{n!\,(n-k)!}\,
\bigl[(-1)^i e_i(\alpha)\bigr]\bigl[(-1)^j e_j(\beta)\bigr]
=\sum_{i+j=k}\frac{(n-i)!\,(n-j)!}{n!\,(n-k)!}\,a_i b_j
= c_k,
\]
which is precisely the coefficient \eqref{eq:coeff} of $p\qfourbn q$. As this holds
for every $k$, the two monic degree-$n$ polynomials in \eqref{eq:walsh} agree.
\end{proof}

\noindent Henceforth we use \eqref{eq:coeff} and \eqref{eq:walsh}
interchangeably. \emph{Walsh's reality theorem} (Lemma~\ref{lem:walsh} below)
guarantees that when $p,q$ are real-rooted, $p\qfourbn q$ is again real-rooted
\cite{Walsh, MSS22}. Thus $\qfourbn$ induces a map
\[
\qfourRootMap:\qfourR^n\times\qfourR^n\to\qfourR^n,
\]
where $\qfourRootMap(\alpha,\beta)$ is the ordered root vector of $p\qfourbn q$.

\smallskip
We now state and prove the reality theorem; it is needed both to make
$\qfourRootMap$ well defined and, more essentially, to supply the hyperbolicity
used in Corollary~\ref{cor:psd}. The result is exactly the content of Walsh's
1922 paper \cite{Walsh}, reproved in modern language in \cite{MSS22}; the proof
we give is the standard one via the Grace--Walsh--Szeg\H{o} coincidence theorem
(itself the central theorem of \cite{Walsh}).

\begin{qfourtheorem}[Grace--Walsh--Szeg\H{o} coincidence theorem
\cite{Walsh}, \cite{RS}, \cite{MSS22}]\label{thm:gws}
Let $\Phi\in\qfourC[y_1,\dots,y_n]$ be \emph{symmetric} (invariant under all
permutations of the $y_i$) and \emph{multiaffine} (of degree at most $1$ in each
variable separately). Let $K\subseteq\qfourC$ be a circular region, i.e.\ a
closed or open disc, half-plane, or the complement of such. If
$\beta_1,\dots,\beta_n\in K$, then there exists a single point $\zeta\in K$ with
\[
\Phi(\beta_1,\dots,\beta_n)\;=\;\Phi(\zeta,\zeta,\dots,\zeta).
\]
\end{qfourtheorem}

\noindent\emph{Intuition.} The hypotheses say $\Phi$ is the unique
\emph{polarization} of a single-variable degree-$\le n$ polynomial: there is a
$\phi(y)$ with $\Phi(y,\dots,y)=\phi(y)$, and $\Phi$ is its symmetric
multiaffine ``spreading out'' over $n$ variables. The theorem says that, as far
as the value $\Phi(\beta_1,\dots,\beta_n)$ is concerned, the $n$ points
$\beta_i$ can always be ``coincided'' into a single point $\zeta$ of the same
circular region. We use it only for the two closed half-planes
$\overline{H_+}=\{\,\Im z\ge0\,\}$ and $\overline{H_-}=\{\,\Im z\le0\,\}$, both of
which contain all of $\qfourR$.

\begin{qfourlemma}[Walsh's reality theorem for $\qfourbn$]\label{lem:walsh}
If $p(x)=\prod_{i=1}^n(x-\alpha_i)$ and $q(x)=\prod_{j=1}^n(x-\beta_j)$ have only
real roots, then every root of $p\qfourbn q$ is real.
\end{qfourlemma}

\begin{proof}
Define the polynomial in $n+1$ variables
\[
G(x;y_1,\dots,y_n)\;:=\;\frac{1}{n!}\sum_{\pi\in S_n}\prod_{i=1}^n\bigl(x-\alpha_i-y_{\pi(i)}\bigr).
\]
For fixed $x$ it is, in the variables $y_1,\dots,y_n$, manifestly symmetric;
moreover it is multiaffine, because in each summand the index $\pi(i)$ runs over
distinct values, so each variable $y_j$ occurs in at most one factor and hence to
degree at most $1$. By construction $G(x;y,\dots,y)=\prod_{i=1}^n(x-\alpha_i-y)$,
so $G$ is the symmetric multiaffine polarization of
$y\mapsto\prod_i(x-\alpha_i-y)$. Comparing with \eqref{eq:walsh},
\begin{equation}\label{eq:G-is-box}
G(x;\beta_1,\dots,\beta_n)\;=\;(p\qfourbn q)(x).
\end{equation}

Let $w\in\qfourC$ be any root of $p\qfourbn q$ and suppose, for contradiction,
that $\Im w\ne0$. Fix this $w$ and put
\[
\Phi(y_1,\dots,y_n)\;:=\;G(w;y_1,\dots,y_n),
\]
a symmetric multiaffine polynomial in $y$ with (in general) complex
coefficients. The Grace--Walsh--Szeg\H{o} theorem holds for complex-coefficient
symmetric multiaffine $\Phi$ --- this is the form in which it is usually stated,
see \cite[Thm.~3.4.1b]{RS} and \cite{MSS22} --- so Theorem~\ref{thm:gws}
applies to $\Phi$. By \eqref{eq:G-is-box} and
$(p\qfourbn q)(w)=0$ we have $\Phi(\beta_1,\dots,\beta_n)=0$.

Apply Theorem~\ref{thm:gws} with the circular region $K=\overline{H_+}$. Since
$\beta_1,\dots,\beta_n$ are real, they lie in $\overline{H_+}$, so there is
$\zeta_+\in\overline{H_+}$ with
\[
\prod_{i=1}^n\bigl(w-\alpha_i-\zeta_+\bigr)=G(w;\zeta_+,\dots,\zeta_+)
=\Phi(\zeta_+,\dots,\zeta_+)=\Phi(\beta_1,\dots,\beta_n)=0 .
\]
Hence $w-\alpha_i-\zeta_+=0$ for some $i$, i.e.\ $\zeta_+=w-\alpha_i$. As
$\alpha_i\in\qfourR$, $\Im\zeta_+=\Im w$, and $\zeta_+\in\overline{H_+}$ forces
$\Im w\ge0$.

Repeating the argument with $K=\overline{H_-}$ (which also contains all the real
$\beta_j$) produces $\zeta_-=w-\alpha_{i'}\in\overline{H_-}$, whence
$\Im w=\Im\zeta_-\le0$. Combining, $\Im w=0$, contradicting $\Im w\ne0$.
Therefore every root of $p\qfourbn q$ is real.
\end{proof}

\noindent\emph{Why a symmetric sum of $n!$ monic polynomials stays real-rooted.}
The phenomenon is not about the sum being a positive combination (it is not, and
sums of real-rooted polynomials are usually \emph{not} real-rooted). The point
is the rigid algebraic structure visible in $G$: the convolution is the value of
a single \emph{multiaffine, symmetric} polynomial at the root vector $\beta$, and
multiaffine symmetric polynomials cannot ``create'' a complex zero out of real
inputs without already having one when all inputs coincide --- this is precisely
the coincidence theorem. Real-rootedness of $p$ keeps the $\alpha_i$ real (so
$\Im\zeta=\Im w$), and real-rootedness of $q$ keeps the $\beta_j$ inside both
half-planes; the two facts pin $\Im w$ to $0$. Write
$m_k(\alpha)=\frac1n\sum_i\alpha_i^k$ and
$\qfourVar(\alpha)=m_2(\alpha)-m_1(\alpha)^2$, and let $\mathbf 1_n$ be the all-ones
vector.

\begin{qfourproposition}[properties of $\qfourbn$]\label{prop:props}
If $\gamma=\qfourRootMap(\alpha,\beta)$ then
\begin{enumerate}\setlength\itemsep{2pt}
\item (additivity) $m_1(\gamma)=m_1(\alpha)+m_1(\beta)$ and
$\qfourVar(\gamma)=\qfourVar(\alpha)+\qfourVar(\beta)$;
\item (translation) $\qfourRootMap(\alpha+t\mathbf 1_n,\beta)=\gamma+t\mathbf 1_n$
and likewise in the second argument.
\end{enumerate}
\end{qfourproposition}

\begin{proof}
Write $\gamma=\qfourRootMap(\alpha,\beta)$ for the (ordered) root vector of
$p\qfourbn q=\sum_{k}c_kx^{n-k}$. By Vi\`ete, the first two power sums of the
$\gamma_\ell$ are recovered from $c_1,c_2$ via the Newton identities
\begin{equation}\label{eq:newton-gamma}
\sum_{\ell}\gamma_\ell=-c_1,\qquad
\sum_{\ell}\gamma_\ell^2=c_1^2-2c_2,
\end{equation}
since $c_1=-e_1(\gamma)$ and $c_2=e_2(\gamma)$, and the power sums satisfy
$p_1=e_1$, $p_2=e_1^2-2e_2$. Thus to prove (i) it suffices to compute $c_1$ and
$c_2$. We do this directly from the symmetrized-product representation
\eqref{eq:walsh}, which is valid by Lemma~\ref{lem:walsh-eq}.

\smallskip\noindent\emph{Setup.}
Put $\gamma_i^\pi:=\alpha_i+\beta_{\pi(i)}$ for $\pi\in S_n$. Expanding the
product in \eqref{eq:walsh},
\[
(p\qfourbn q)(x)=\frac1{n!}\sum_{\pi\in S_n}\prod_{i=1}^n\bigl(x-\gamma_i^\pi\bigr)
=\sum_{k=0}^n(-1)^k\Bigl(\frac1{n!}\sum_{\pi\in S_n}e_k(\gamma^\pi)\Bigr)x^{n-k},
\]
so $c_k=(-1)^k\,\overline{e_k}$ where $\overline{e_k}:=\frac1{n!}\sum_{\pi}e_k(\gamma^\pi)$.
We compute $\overline{e_1}$ and $\overline{e_2}$. Throughout we abbreviate
$e_1(\alpha)=\sum_i\alpha_i$, $e_2(\alpha)=\sum_{i<j}\alpha_i\alpha_j$, and
likewise for $\beta$.

\smallskip\noindent\emph{First coefficient.}
For each $\pi$,
$e_1(\gamma^\pi)=\sum_{i=1}^n(\alpha_i+\beta_{\pi(i)})
=\sum_i\alpha_i+\sum_i\beta_{\pi(i)}=e_1(\alpha)+e_1(\beta)$,
since $\pi$ permutes the $\beta$'s. The right side is independent of $\pi$, so
averaging gives
\begin{equation}\label{eq:c1}
\overline{e_1}=e_1(\alpha)+e_1(\beta),\qquad
c_1=-\bigl(e_1(\alpha)+e_1(\beta)\bigr).
\end{equation}
By \eqref{eq:newton-gamma}, $\sum_\ell\gamma_\ell=e_1(\alpha)+e_1(\beta)=\sum_i\alpha_i+\sum_j\beta_j$,
and dividing by $n$ yields $m_1(\gamma)=m_1(\alpha)+m_1(\beta)$.

\smallskip\noindent\emph{Second coefficient.}
Fix $\pi$ and expand $e_2$ of the sums:
\[
e_2(\gamma^\pi)=\sum_{i<j}(\alpha_i+\beta_{\pi(i)})(\alpha_j+\beta_{\pi(j)})
=\underbrace{\sum_{i<j}\alpha_i\alpha_j}_{=\,e_2(\alpha)}
+\underbrace{\sum_{i<j}\beta_{\pi(i)}\beta_{\pi(j)}}_{(\ast)}
+\underbrace{\sum_{i<j}\bigl(\alpha_i\beta_{\pi(j)}+\alpha_j\beta_{\pi(i)}\bigr)}_{(\ast\ast)} .
\]
The first term is already $\pi$-independent and equals $e_2(\alpha)$. For the
second term, $\{\pi(i):i\}=[n]$, and as $(i,j)$ ranges over unordered pairs so
does $(\pi(i),\pi(j))$; hence $(\ast)=\sum_{i<j}\beta_{\pi(i)}\beta_{\pi(j)}
=\sum_{k<\ell}\beta_k\beta_\ell=e_2(\beta)$, again independent of $\pi$. Therefore
\begin{equation}\label{eq:e2-avg-split}
\overline{e_2}=e_2(\alpha)+e_2(\beta)+\frac1{n!}\sum_{\pi\in S_n}(\ast\ast).
\end{equation}

It remains to average the cross term $(\ast\ast)$. We use the elementary
single-index average: for any fixed index $i\in[n]$,
\begin{equation}\label{eq:single-index}
\frac1{n!}\sum_{\pi\in S_n}\beta_{\pi(i)}
=\frac1{n!}\,(n-1)!\sum_{k=1}^n\beta_k
=\frac1n\,e_1(\beta),
\end{equation}
because for each value $k$ there are exactly $(n-1)!$ permutations with
$\pi(i)=k$. Now interchange the (finite) sums and apply \eqref{eq:single-index}
to each summand:
\[
\frac1{n!}\sum_{\pi}(\ast\ast)
=\sum_{i<j}\Bigl(\alpha_i\cdot\tfrac1{n!}\sum_\pi\beta_{\pi(j)}
+\alpha_j\cdot\tfrac1{n!}\sum_\pi\beta_{\pi(i)}\Bigr)
=\sum_{i<j}\bigl(\alpha_i+\alpha_j\bigr)\cdot\frac{e_1(\beta)}{n}.
\]
(Here we used that $\tfrac1{n!}\sum_\pi\beta_{\pi(i)}=\tfrac1{n!}\sum_\pi\beta_{\pi(j)}=\tfrac1n e_1(\beta)$
by \eqref{eq:single-index}, valid since $i\ne j$.) Finally,
\[
\sum_{i<j}(\alpha_i+\alpha_j)=(n-1)\sum_{i}\alpha_i=(n-1)\,e_1(\alpha),
\]
because each index $i$ appears in exactly $n-1$ unordered pairs. Hence
\begin{equation}\label{eq:cross-avg}
\frac1{n!}\sum_{\pi}(\ast\ast)=\frac{n-1}{n}\,e_1(\alpha)\,e_1(\beta).
\end{equation}
Substituting \eqref{eq:cross-avg} into \eqref{eq:e2-avg-split} and using
$c_2=(-1)^2\overline{e_2}=\overline{e_2}$,
\begin{equation}\label{eq:c2}
c_2=e_2(\alpha)+e_2(\beta)+\frac{n-1}{n}\,e_1(\alpha)\,e_1(\beta).
\end{equation}

\smallskip\noindent\emph{Variance additivity.}
Abbreviate $A_1:=e_1(\alpha),\,A_2:=e_2(\alpha)$ and $B_1:=e_1(\beta),\,B_2:=e_2(\beta)$.
For any degree-$n$ vector $v$ one has, by Newton's identities,
\[
n\,m_1(v)=e_1(v),\qquad n\,m_2(v)=p_2(v)=e_1(v)^2-2e_2(v),
\]
so
\[
\qfourVar(v)=m_2(v)-m_1(v)^2=\frac{e_1(v)^2-2e_2(v)}{n}-\frac{e_1(v)^2}{n^2}
=\frac{(n-1)e_1(v)^2-2n\,e_2(v)}{n^2}.
\]
In particular
\[
\qfourVar(\alpha)=\frac{(n-1)A_1^2-2nA_2}{n^2},\qquad
\qfourVar(\beta)=\frac{(n-1)B_1^2-2nB_2}{n^2}.
\]
For $\gamma$, by \eqref{eq:c1} and \eqref{eq:c2} we have
$e_1(\gamma)=-c_1=A_1+B_1$ and $e_2(\gamma)=c_2=A_2+B_2+\tfrac{n-1}{n}A_1B_1$, so
\begin{align*}
\qfourVar(\gamma)
&=\frac{(n-1)(A_1+B_1)^2-2n\bigl(A_2+B_2+\tfrac{n-1}{n}A_1B_1\bigr)}{n^2}\\
&=\frac{(n-1)\bigl(A_1^2+2A_1B_1+B_1^2\bigr)-2nA_2-2nB_2-2(n-1)A_1B_1}{n^2}\\
&=\frac{(n-1)A_1^2+(n-1)B_1^2-2nA_2-2nB_2
+\bigl[2(n-1)-2(n-1)\bigr]A_1B_1}{n^2}\\
&=\frac{(n-1)A_1^2-2nA_2}{n^2}+\frac{(n-1)B_1^2-2nB_2}{n^2}
=\qfourVar(\alpha)+\qfourVar(\beta).
\end{align*}
The cross term $A_1B_1$ cancels exactly, which is the whole point of the
$\tfrac{n-1}{n}$ factor in \eqref{eq:c2}. This proves the variance statement and
completes the proof of (i).

\smallskip\noindent
(ii) Translating $\alpha\mapsto\alpha+t\mathbf 1_n$ replaces each
$\gamma_i^\pi=\alpha_i+\beta_{\pi(i)}$ by $\gamma_i^\pi+t$, so every factor
$x-\gamma_i^\pi$ in \eqref{eq:walsh} becomes $(x-t)-\gamma_i^\pi$; thus
$(\,(p\text{ shifted})\qfourbn q)(x)=(p\qfourbn q)(x-t)$, whose root vector is
$\gamma+t\mathbf 1_n$. The same argument applies in the second argument, giving
$\qfourRootMap(\alpha+t\mathbf 1_n,\beta)=\gamma+t\mathbf 1_n$.
\end{proof}

\section*{Score vectors via the heat flow}

The link between $\qfourInfo$ and the dynamics of roots is the following, due to
Shlyakhtenko.

\begin{qfourlemma}[score vector as a velocity]\label{lem:heat}
Assume $p$ has simple roots. Let $p_t(x):=\exp\!\bigl(-\tfrac{t}{2}\partial_x^2\bigr)p(x)$
and let $\alpha(t)$ be the ordered root vector of $p_t$. Then
$\alpha_i'(0)=\sum_{j\ne i}\frac{1}{\alpha_i-\alpha_j}$; that is,
$\alpha'(0)=\qfourScore(\alpha)$.
\end{qfourlemma}

\begin{proof}
Roots stay simple near $t=0$. By definition, $p_t(\alpha_i(t))=0$. The heat operator
expands as $p_t(x)=\exp(-\tfrac{t}{2}\partial_x^2)p(x)=p(x)-\tfrac{t}{2}p''(x)+O(t^2)$.
Thus at $x=\alpha_i(t)$,
\[
p(\alpha_i(t))-\tfrac{t}{2}p''(\alpha_i(t))+O(t^2)=0.
\]
Differentiating with respect to $t$ at $t=0$:
\[
p'(\alpha_i)\alpha_i'(0)-\tfrac{1}{2}p''(\alpha_i)=0.
\]
Thus $\alpha_i'(0)=\tfrac{1}{2}\,p''(\alpha_i)/p'(\alpha_i)$.

Now, for the polynomial $p(x)=\prod_j(x-\alpha_j)$, we use the factorization
$p(x)=(x-\alpha_i)q(x)$ where $q(x)=\prod_{j\ne i}(x-\alpha_j)$. Then
$p'(\alpha_i)=q(\alpha_i)$ and $p''(\alpha_i)=2q'(\alpha_i)$ (from $p'(x)=q(x)+(x-\alpha_i)q'(x)$
and $p''(x)=2q'(x)+(x-\alpha_i)q''(x)$). Since $q'(x)/q(x)=\sum_{j\ne i}\tfrac{1}{x-\alpha_j}$,
evaluating at $x=\alpha_i$ gives
\[
\frac{p''(\alpha_i)}{p'(\alpha_i)}=\frac{2q'(\alpha_i)}{q(\alpha_i)}=2\sum_{j\ne i}\frac{1}{\alpha_i-\alpha_j}.
\]
Therefore $\alpha_i'(0)=\tfrac{1}{2}\cdot 2\sum_{j\ne i}\tfrac{1}{\alpha_i-\alpha_j}=\sum_{j\ne i}\tfrac{1}{\alpha_i-\alpha_j}$.
\end{proof}

The essential point is that $\exp(-\tfrac t2\partial_x^2)$ \emph{commutes with}
$\qfourbn$ in the precise sense
$(e^{-at/2\partial_x^2}p)\qfourbn(e^{-bt/2\partial_x^2}q)=e^{-(a+b)t/2\partial_x^2}(p\qfourbn q)$
(proved in Corollary~\ref{cor:heat-commutes} below by the same coefficient
argument as Lemma~\ref{lem:simple}(iii)). We exploit this twice.

\section*{Reduction to simple roots}

We first record an elementary but crucial quantitative estimate: a small root
gap forces $\qfourInfo$ to be large. The naive bound
$\qfourInfo(\alpha)\ge\bigl(1/(\alpha_i-\alpha_j)\bigr)^2$ is \emph{not} valid,
because the $i$-th score $\qfourScore(\alpha)_i=\sum_{l\ne i}1/(\alpha_i-\alpha_l)$
is a \emph{sum} in which the large term $1/(\alpha_i-\alpha_j)$ may be partly
cancelled by other terms (e.g.\ for the symmetric triple
$\alpha=(-\delta,0,\delta)$ the middle score is exactly $0$). The correct
statement is obtained by a scaling argument applied to whole clusters of nearly
coincident roots.

\begin{qfourlemma}[gap blow-up]\label{lem:gapblowup}
Let $\xi^{(k)}\in\qfourR^n$ be a sequence of vectors with \emph{distinct} entries
converging to a limit $\xi\in\qfourR^n$ that has a repeated entry. Then
$\qfourInfo(\xi^{(k)})=\|\qfourScore(\xi^{(k)})\|^2\to+\infty$. More precisely,
if a class of mutually coincident limit coordinates has size $m\ge2$ and the
corresponding cluster of $\xi^{(k)}$-coordinates has diameter $d_k\to0$, then
$\qfourInfo(\xi^{(k)})\ge \rho_m/(2d_k^2)-O(1)$ for a constant $\rho_m\ge1$.
\end{qfourlemma}

\begin{proof}
Partition $\{1,\dots,n\}$ into the level sets of the limit: $i\sim j$ iff
$\xi_i=\xi_j$. Let $G$ be a class with $|G|=m\ge 2$ (one exists since $\xi$ has a
repeated entry) and common limit value $v$, so $\xi^{(k)}_i\to v$ for all $i\in G$,
while for $l\notin G$ the limit $\xi_l\ne v$.

\smallskip\noindent\emph{Internal vs.\ external part of the score.}
For $i\in G$ write
\[
\qfourScore(\xi^{(k)})_i
=\underbrace{\sum_{l\in G,\,l\ne i}\frac1{\xi^{(k)}_i-\xi^{(k)}_l}}_{=:I^{(k)}_i\ \text{(internal)}}
\;+\;
\underbrace{\sum_{l\notin G}\frac1{\xi^{(k)}_i-\xi^{(k)}_l}}_{=:E^{(k)}_i\ \text{(external)}}.
\]
Since $\xi^{(k)}_i\to v$ and $\xi^{(k)}_l\to\xi_l\ne v$ for $l\notin G$, each
external denominator converges to the nonzero number $v-\xi_l$; hence the external
parts are bounded: there is $C<\infty$ with $|E^{(k)}_i|\le C$ for all $i\in G$ and
all large $k$.

\smallskip\noindent\emph{The internal parts blow up, in the $\ell^2$ sense.}
Let $\eta^{(k)}:=(\xi^{(k)}_i)_{i\in G}\in\qfourR^m$ be the sub-vector of cluster
coordinates, with distinct entries, and let
$d_k:=\max_{i,j\in G}|\xi^{(k)}_i-\xi^{(k)}_j|>0$ be its diameter; then $d_k\to 0$
because all entries tend to $v$. Note $(I^{(k)}_i)_{i\in G}=\qfourScore(\eta^{(k)})$
is exactly the score vector of the $m$-point configuration $\eta^{(k)}$, so
\[
\sum_{i\in G}\bigl(I^{(k)}_i\bigr)^2=\qfourInfo(\eta^{(k)})
\quad(\text{Fisher information of the cluster, in dimension }m).
\]
Normalise: write $\eta^{(k)}=\mu_k\mathbf 1+d_k\,\zeta^{(k)}$ where
$\mu_k=\min_{i\in G}\xi^{(k)}_i$ and $\zeta^{(k)}\in[0,1]^m$ has entries spanning
$[0,1]$ (so $\zeta^{(k)}$ has diameter exactly $1$, in particular distinct entries
and not all equal). The score is translation invariant and scales by $1/d_k$ under
the dilation $\eta\mapsto\mu+d\,\zeta$ (each difference scales by $d_k$), so
\[
\qfourInfo(\eta^{(k)})=\frac{1}{d_k^2}\,\qfourInfo(\zeta^{(k)}).
\]
The normalised configurations $\zeta^{(k)}$ lie in the compact set
$K_m=\{\zeta\in[0,1]^m:\ \min_i\zeta_i=0,\ \max_i\zeta_i=1\}$. Define
$\rho_m:=\inf_{\zeta\in K_m'}\qfourInfo(\zeta)$, where
$K_m'=\{\zeta\in K_m:\ \text{entries of }\zeta\text{ distinct}\}$. We claim
$\rho_m>0$. Indeed, for any \emph{distinct} $\zeta\in K_m'$, order the entries
$\zeta_{(1)}<\dots<\zeta_{(m)}$; the largest one has score
$\qfourScore(\zeta)_{(m)}=\sum_{l<m}1/(\zeta_{(m)}-\zeta_{(l)})>0$ (every summand
positive), and since $\zeta_{(m)}-\zeta_{(l)}\le\mathrm{diam}=1$ each summand is
$\ge 1$, so $\qfourScore(\zeta)_{(m)}\ge m-1\ge 1$, whence
$\qfourInfo(\zeta)\ge(m-1)^2\ge 1$. Thus $\rho_m\ge 1>0$. Consequently
\[
\sum_{i\in G}\bigl(I^{(k)}_i\bigr)^2=\qfourInfo(\eta^{(k)})
=\frac{\qfourInfo(\zeta^{(k)})}{d_k^2}\ \ge\ \frac{\rho_m}{d_k^2}
\ \xrightarrow[k\to\infty]{}\ +\infty .
\]

\smallskip\noindent\emph{Conclusion.}
Using $(a+b)^2\ge\tfrac12 a^2-b^2$ coordinatewise on the cluster $G$,
\[
\qfourInfo(\xi^{(k)})\ge\sum_{i\in G}\qfourScore(\xi^{(k)})_i^2
=\sum_{i\in G}\bigl(I^{(k)}_i+E^{(k)}_i\bigr)^2
\ge\tfrac12\sum_{i\in G}\bigl(I^{(k)}_i\bigr)^2-\sum_{i\in G}\bigl(E^{(k)}_i\bigr)^2
\ge\frac{\rho_m}{2\,d_k^2}-mC^2 .
\]
As $d_k\to0$ the right-hand side tends to $+\infty$, so
$\qfourInfo(\xi^{(k)})\to+\infty$, with the quantitative bound
$\qfourInfo(\xi^{(k)})\ge\rho_m/(2d_k^2)-mC^2$ recorded in the statement.
\end{proof}


\begin{qfourlemma}\label{lem:simple}
Let $T_\epsilon:=(1-\epsilon\,d/dx)^n$. For monic real-rooted $p$ of degree $n$:
\begin{enumerate}\setlength\itemsep{2pt}
\item $T_\epsilon p$ is monic real-rooted with simple roots;
\item $\qfourInfo$ is lower semicontinuous in the roots: if $\alpha_k\to\alpha$,
then $\liminf_{k\to\infty}\qfourInfo(p_k)\ge\qfourInfo(p)$;
\item $(T_\epsilon p)\qfourbn(T_\epsilon q)=T_\epsilon^2(p\qfourbn q)$.
\end{enumerate}
\end{qfourlemma}

\begin{proof}
(i) is Nuij's theorem \cite{Nuij}.

(ii) Let $\alpha_k=(\alpha_{k,1},\dots,\alpha_{k,n})\to\alpha=(\alpha_1,\dots,\alpha_n)$.
If $\alpha$ has simple roots, then the score function
$\qfourScore(\alpha)_i=\sum_{j\ne i}\frac{1}{\alpha_i-\alpha_j}$ is continuous in
$\alpha$ (in a neighbourhood of simple-root configurations), so
$\qfourInfo(p_k)=\|\qfourScore(\alpha_k)\|^2\to\|\qfourScore(\alpha)\|^2=\qfourInfo(p)$.

If $\alpha$ has a repeated root, then $\qfourInfo(p)=\infty$. We must show
$\qfourInfo(p_k)\to\infty$. If some $p_k$ also has a repeated root we interpret
$\qfourInfo(p_k)=\infty$ and may discard those terms; on the cofinal set of indices
for which $\alpha_k$ has distinct entries, the gap blow-up estimate
Lemma~\ref{lem:gapblowup} (applied to the sequence $\alpha_k\to\alpha$, whose limit
$\alpha$ has a repeated entry) gives $\qfourInfo(p_k)=\|\qfourScore(\alpha_k)\|^2\to+\infty$.
In all cases $\qfourInfo(p_k)\to+\infty$, so
$\liminf_{k\to\infty}\qfourInfo(p_k)=\infty=\qfourInfo(p)$.
(The naive single-term bound $\qfourInfo(p_k)\ge(1/(\alpha_{k,i}-\alpha_{k,j}))^2$ is
\emph{invalid}, since the $i$-th score is a sum that may exhibit cancellation;
Lemma~\ref{lem:gapblowup} circumvents this by an $\ell^2$ cluster-scaling argument.)

In both cases, $\liminf_{k\to\infty}\qfourInfo(p_k)\ge\qfourInfo(p)$, proving LSC.

(iii) We prove this directly from the coefficient formula, with no separate
appeal to MSS22.  Write $p(x)=\sum_{k=0}^n a_k x^{n-k}$ and
$q(x)=\sum_{k=0}^n b_k x^{n-k}$, and set
$c_k = \sum_{i+j=k}\frac{(n-i)!(n-j)!}{n!(n-k)!}a_i b_j$ for the
coefficients of $p\qfourbn q$.

\smallskip\noindent\textit{Action of $T_\epsilon$ on coefficients.}
Expanding $T_\epsilon=(1-\epsilon\,d/dx)^n=\sum_{m=0}^n\binom{n}{m}(-\epsilon)^m(d/dx)^m$
and using $(d/dx)^m x^{n-j}=\tfrac{(n-j)!}{(n-j-m)!}x^{n-j-m}$, the
coefficient of $x^{n-k}$ in $T_\epsilon r$ is
\[
[T_\epsilon r]_k
\;=\;\sum_{m=0}^k \binom{n}{m}(-\epsilon)^m\frac{(n{-}k{+}m)!}{(n{-}k)!}\,r_{k-m}.
\]

\smallskip\noindent\textit{Coefficient-level computation of the LHS.}
We adopt throughout the convention that $a_i=b_j=0$ whenever $i<0$ or $i>n$
(resp.\ $j<0$ or $j>n$), and that $\binom{n}{m}=0$ for $m<0$ or $m>n$; with
these conventions all the sums below may be written over unrestricted
nonnegative indices, the out-of-range terms vanishing automatically. The formula
for $\qfourbn$ is bilinear in the coefficient sequences of its arguments, so
\[
[(T_\epsilon p)\qfourbn(T_\epsilon q)]_k
= \sum_{i+j=k}\frac{(n-i)!(n-j)!}{n!(n-k)!}\,[T_\epsilon p]_i\,[T_\epsilon q]_j ,
\]
the sum being over $0\le i,j\le n$ with $i+j=k$ (so $0\le i,j\le k\le n$, and the
factorials $(n-i)!,(n-j)!$ are all well defined). Substituting the formula
\[
[T_\epsilon r]_i=\sum_{m_1=0}^{i}\binom{n}{m_1}(-\epsilon)^{m_1}\frac{(n-i+m_1)!}{(n-i)!}\,r_{i-m_1}
\]
for $[T_\epsilon p]_i$ (with index $m_1$) and the analogous one for
$[T_\epsilon q]_j$ (with index $m_2$), we obtain the quadruple sum
\begin{equation}\label{eq:lhs-quad}
[(T_\epsilon p)\qfourbn(T_\epsilon q)]_k
=\frac{1}{n!(n-k)!}\!\!\sum_{\substack{i+j=k\\ 0\le i,j}}\;
\sum_{\substack{0\le m_1\le i\\ 0\le m_2\le j}}
\binom{n}{m_1}\binom{n}{m_2}(-\epsilon)^{m_1+m_2}
(n{-}i{+}m_1)!\,(n{-}j{+}m_2)!\,a_{i-m_1}b_{j-m_2},
\end{equation}
where we used that $\frac{(n-i)!(n-j)!}{n!(n-k)!}\cdot\frac{(n-i+m_1)!}{(n-i)!}\cdot\frac{(n-j+m_2)!}{(n-j)!}
=\frac{(n-i+m_1)!(n-j+m_2)!}{n!(n-k)!}$.

\smallskip\noindent\emph{Change of variables.}
Put $i'=i-m_1$, $j'=j-m_2$ and $m=m_1+m_2$. As $(i,j,m_1,m_2)$ ranges over the
index set of \eqref{eq:lhs-quad}, the tuple $(i',j',m_1,m_2)$ ranges over
\[
\Omega_k=\bigl\{(i',j',m_1,m_2):\ i'\ge0,\ j'\ge0,\ m_1\ge0,\ m_2\ge0,\ i'+j'+m_1+m_2=k\bigr\},
\]
this being a bijection with inverse $i=i'+m_1$, $j=j'+m_2$ (the constraints
$m_1\le i$, $m_2\le j$ are exactly $i'\ge0$, $j'\ge0$, and $i+j=k$ becomes
$i'+j'+m_1+m_2=k$). Under it, $n-i+m_1=n-i'$ and $n-j+m_2=n-j'$, so
\eqref{eq:lhs-quad} becomes
\[
[(T_\epsilon p)\qfourbn(T_\epsilon q)]_k
=\frac{1}{n!(n-k)!}\sum_{(i',j',m_1,m_2)\in\Omega_k}
\binom{n}{m_1}\binom{n}{m_2}(-\epsilon)^{m_1+m_2}(n-i')!(n-j')!\,a_{i'}b_{j'}.
\]

\smallskip\noindent\emph{Grouping by $m=m_1+m_2$.}
For fixed $m\in\{0,\dots,k\}$, the slice of $\Omega_k$ with $m_1+m_2=m$ factors as
a free product: $(m_1,m_2)$ ranges over all pairs with $m_1+m_2=m$, $m_1,m_2\ge0$,
\emph{independently} of $(i',j')$, which ranges over all pairs with $i'+j'=k-m$,
$i',j'\ge0$. (The two blocks of constraints share no variables and only fix the
two partial sums $m_1+m_2=m$ and $i'+j'=k-m$ separately.) Crucially, the
$(-\epsilon)^{m_1+m_2}=(-\epsilon)^m$ and $a_{i'}b_{j'}(n-i')!(n-j')!$ factors
depend only on $m$ and on $(i',j')$, not on the particular split $(m_1,m_2)$;
hence the $(m_1,m_2)$-sum and the $(i',j')$-sum separate:
\begin{equation}\label{eq:lhs-vandermonde}
[(T_\epsilon p)\qfourbn(T_\epsilon q)]_k
= \frac{1}{n!(n-k)!}\sum_{m=0}^k(-\epsilon)^m
   \Bigl(\sum_{m_1=0}^m\binom{n}{m_1}\binom{n}{m-m_1}\Bigr)
   \sum_{i'+j'=k-m}(n-i')!(n-j')!\,a_{i'}b_{j'}.
\end{equation}
Three remarks justify every range above. (a) Out-of-range terms vanish, not are
discarded: if $m_1>n$ then $\binom{n}{m_1}=0$, so extending the inner $m_1$-sum
to all $0\le m_1\le m$ (rather than $m_1\le\min(m,n)$) changes nothing; likewise
$\binom{n}{m-m_1}=0$ when $m-m_1>n$. (b) The $(i',j')$-sum runs over $i'+j'=k-m$
with $i',j'\ge0$, and the upper bounds $i'\le n$, $j'\le n$ are automatic because
$a_{i'}=0$ for $i'>n$ and $b_{j'}=0$ for $j'>n$ (and indeed $i'\le k-m\le k\le n$
already). Since the limit $i'+j'=k-m$ and the requirement $i',j'\ge0$ do not
involve $(m_1,m_2)$, the inner $(i',j')$-sum is the same for every split, which is
precisely what licenses pulling it outside the $(m_1,m_2)$-sum. (c) The outer
index $m$ runs only over $0\le m\le k$: for $m>k$ the slice of $\Omega_k$ is
empty (it would force $i'+j'=k-m<0$). By Vandermonde's convolution identity,
$\sum_{m_1=0}^m\binom{n}{m_1}\binom{n}{m-m_1}=\binom{2n}{m}$.

\smallskip\noindent\textit{Coefficient-level computation of the RHS.}
Since $p\qfourbn q$ has degree $n$, applying
$T_\epsilon^2=(1-\epsilon\,d/dx)^{2n}$ gives
\[
[T_\epsilon^2(p\qfourbn q)]_k
= \sum_{m=0}^k\binom{2n}{m}(-\epsilon)^m\frac{(n{-}k{+}m)!}{(n{-}k)!}\,c_{k-m}.
\]
Substituting $c_{k-m}=\sum_{i'+j'=k-m}\frac{(n-i')!(n-j')!}{n!(n-k+m)!}a_{i'}b_{j'}$
and cancelling the factor $(n{-}k{+}m)!$ appearing in both $\frac{(n{-}k{+}m)!}{(n{-}k)!}$
and the denominator of $c_{k-m}$, leaving $\frac{1}{(n{-}k)!}$:
\[
[T_\epsilon^2(p\qfourbn q)]_k
= \frac{1}{n!(n-k)!}\sum_{m=0}^k\binom{2n}{m}(-\epsilon)^m
  \sum_{i'+j'=k-m}(n-i')!(n-j')!\,a_{i'}b_{j'}.
\]
This is identical to~\eqref{eq:lhs-vandermonde} after applying Vandermonde,
so $[(T_\epsilon p)\qfourbn(T_\epsilon q)]_k=[T_\epsilon^2(p\qfourbn q)]_k$ for
every $k\in\{0,\dots,n\}$.

\smallskip\noindent\textit{Why $T_\epsilon^2$ and not a different power.}
Each of the two arguments of $\qfourbn$ contributes a binomial weight
$\binom{n}{m_i}(-\epsilon)^{m_i}$; their convolution over $m_1+m_2=m$
equals $\binom{2n}{m}(-\epsilon)^m$ by Vandermonde — exactly the weight
appearing in $(1-\epsilon\,d/dx)^{2n}=T_\epsilon^2$.
(The same argument applies to the heat operator; see
Corollary~\ref{cor:heat-commutes} immediately below.)
\end{proof}

\begin{qfourcorollary}[heat-flow commutativity]\label{cor:heat-commutes}
For all $a,b\ge0$ and $s\ge0$,
\[
\bigl(e^{-as/2\,\partial_x^2}p\bigr)\qfourbn\bigl(e^{-bs/2\,\partial_x^2}q\bigr)
\;=\;e^{-(a+b)s/2\,\partial_x^2}\bigl(p\qfourbn q\bigr).
\]
\end{qfourcorollary}

\begin{proof}
The argument mirrors part~(iii) above, with exponential weights in place of
binomial weights.  Write $p(x)=\sum_k a_kx^{n-k}$, $q(x)=\sum_k b_kx^{n-k}$.
The $k$-th coefficient of $e^{-s/2\,\partial_x^2}r$ is
\[
[e^{-s/2\,\partial_x^2}r]_k
\;=\;\sum_{m=0}^{\lfloor k/2\rfloor}\frac{(-s/2)^m}{m!}
 \frac{(n{-}k{+}2m)!}{(n{-}k)!}\,r_{k-2m},
\]
obtained from $(d/dx)^{2m}[r_{k-2m}x^{n-k+2m}]
=r_{k-2m}\frac{(n-k+2m)!}{(n-k)!}x^{n-k}$.
Substituting into the bilinear formula for $\qfourbn$ and
changing variables $i'=i-2m_1$, $j'=j-2m_2$, $\ell=m_1+m_2$
(so $i'+j'=k-2\ell$, $n-i+2m_1=n-i'$, $n-j+2m_2=n-j'$; the same
bijection-and-grouping argument as in Lemma~\ref{lem:simple}(iii) applies, with
the series in $m_1,m_2$ now terminating because $r_{i-2m_1}=0$ once
$i-2m_1<0$ and $a_{i'}=b_{j'}=0$ outside $0\le i',j'\le n$) gives
\begin{align*}
\bigl[(e^{-as/2\partial_x^2}p)\qfourbn(e^{-bs/2\partial_x^2}q)\bigr]_k
&\;=\;\frac{1}{n!(n{-}k)!}\sum_{\ell\ge0}(-s/2)^\ell
 \!\!\!\sum_{m_1+m_2=\ell}\!\!\frac{a^{m_1}b^{m_2}}{m_1!\,m_2!}
 \sum_{i'+j'=k-2\ell}(n{-}i')!(n{-}j')!\,a_{i'}b_{j'}.
\end{align*}
The exponential binomial identity (expanding $e^{au}e^{bu}=e^{(a+b)u}$
and extracting the $u^\ell$ coefficient) gives
$\sum_{m_1+m_2=\ell}\frac{a^{m_1}b^{m_2}}{m_1!\,m_2!}=\frac{(a+b)^\ell}{\ell!}$,
so the above becomes
\[
\frac{1}{n!(n{-}k)!}\sum_{\ell\ge0}\frac{(-(a+b)s/2)^\ell}{\ell!}
\sum_{i'+j'=k-2\ell}(n{-}i')!(n{-}j')!\,a_{i'}b_{j'}.
\]
For the right-hand side, $c_{k-2\ell}=\sum_{i'+j'=k-2\ell}
\frac{(n-i')!(n-j')!}{n!(n-k+2\ell)!}a_{i'}b_{j'}$, so
\[
[e^{-(a+b)s/2\partial_x^2}(p\qfourbn q)]_k
=\sum_\ell\frac{(-(a+b)s/2)^\ell}{\ell!}\frac{(n{-}k{+}2\ell)!}{(n{-}k)!}
\cdot\frac{\sum_{i'+j'=k-2\ell}(n{-}i')!(n{-}j')!\,a_{i'}b_{j'}}{n!(n{-}k{+}2\ell)!},
\]
and the $(n{-}k{+}2\ell)!$ factors cancel, leaving an expression identical
to the left-hand side.  Hence both sides agree at every
$k\in\{0,\dots,n\}$.
\end{proof}

\noindent\textbf{Completion of the reduction.}
We first record the elementary topological facts we need, then carry out the
limit.

\medskip
\noindent\textit{Continuity of ordered roots in the coefficients.}
For a monic polynomial $r(x)=x^n+\sum_{k=1}^n r_k x^{n-k}$ identify $r$ with its
coefficient vector $(r_1,\dots,r_n)\in\qfourR^n$. We use the standard fact that
the roots of a monic polynomial depend continuously on its coefficients: if a
sequence of monic polynomials $r^{(m)}\to r$ coefficientwise, then for a suitable
labelling the (multiset of) roots of $r^{(m)}$ converges to the roots of $r$
\cite[Ch.~II, App.~A]{Whitney} (a consequence of Rouch\'e's theorem). When all
polynomials in question are \emph{real-rooted} of the same degree $n$ and we use
the canonical ordering $\lambda_1\le\dots\le\lambda_n$ of the real roots, this
labelling is the ordered one; thus the map sending a real-rooted monic
coefficient vector to its ordered root vector is continuous. In particular, if
$r^{(m)}\to r$ coefficientwise through real-rooted monic polynomials of degree
$n$, then the ordered root vectors satisfy $\alpha^{(m)}\to\alpha$ in $\qfourR^n$.

\medskip
\noindent\textit{The $\epsilon$-inequality.}
By Lemma~\ref{lem:simple}(i), $T_\epsilon p$, $T_\epsilon q$ and
$T_\epsilon^2(p\qfourbn q)$ are monic real-rooted with \emph{simple} roots for
every $\epsilon>0$. By Lemma~\ref{lem:simple}(iii),
$(T_\epsilon p)\qfourbn(T_\epsilon q)=T_\epsilon^2(p\qfourbn q)$; this is precisely
how applying $T_\epsilon$ to each factor relates to applying $T_\epsilon^2$ to the
convolution. Hence the simple-root case of the theorem (proved in §§1--4 below
for the simple-rooted triple $T_\epsilon p,\,T_\epsilon q,\,(T_\epsilon p)\qfourbn(T_\epsilon q)$)
yields, for every $\epsilon>0$,
\begin{equation}\label{eq:eps-ineq}
\frac{1}{\qfourInfo\bigl(T_\epsilon^2(p\qfourbn q)\bigr)}\;\ge\;
\frac{1}{\qfourInfo(T_\epsilon p)}+\frac{1}{\qfourInfo(T_\epsilon q)}.
\end{equation}

We now pass to the limit $\epsilon\to 0$ in three steps.

\medskip
\noindent\textit{Step A: convergence of approximants (coefficients and roots).}
For any monic polynomial $r$ of degree $n$, the binomial expansion
\[
T_\epsilon r \;=\; (1-\epsilon\,d/dx)^n r \;=\; \sum_{k=0}^{n}\binom{n}{k}(-\epsilon)^k r^{(k)}
\]
shows that each coefficient of $T_\epsilon r$ is a polynomial in $\epsilon$ whose
$\epsilon^0$ term is the corresponding coefficient of $r$; hence
$T_\epsilon r\to r$ coefficientwise as $\epsilon\to 0$. Applying this to
$r\in\{p,\,q,\,p\qfourbn q\}$ (and using $T_\epsilon^2=(1-\epsilon\,d/dx)^{2n}$ for
the last, which likewise tends coefficientwise to $p\qfourbn q$) gives
\[
T_\epsilon p\to p,\qquad T_\epsilon q\to q,\qquad T_\epsilon^2(p\qfourbn q)\to p\qfourbn q
\quad\text{(coefficientwise).}
\]
All polynomials involved are real-rooted monic of degree $n$, so by the
continuity of ordered roots above, the corresponding ordered root vectors
converge:
\[
\alpha^\epsilon:=\mathrm{roots}(T_\epsilon p)\to\alpha,\quad
\beta^\epsilon:=\mathrm{roots}(T_\epsilon q)\to\beta,\quad
\gamma^\epsilon:=\mathrm{roots}\bigl(T_\epsilon^2(p\qfourbn q)\bigr)\to\gamma,
\]
where $\alpha,\beta,\gamma$ are the ordered root vectors of $p,q,p\qfourbn q$.

\medskip
\noindent\textit{Step B: convergence of the right-hand side.}
We prove the following lemma, which is the precise convergence statement needed.

\begin{qfourlemma}[convergence of $1/\qfourInfo$ along the Nuij approximation]\label{lem:Bconv}
For every monic real-rooted $p$ of degree $n$,
\[
\lim_{\epsilon\to 0}\frac{1}{\qfourInfo(T_\epsilon p)}=\frac{1}{\qfourInfo(p)},
\qquad\text{with the convention }1/\infty=0 .
\]
\end{qfourlemma}

\begin{proof}
By Step~A the ordered root vectors converge, $\alpha^\epsilon\to\alpha$, where
$\alpha^\epsilon=\mathrm{roots}(T_\epsilon p)$ and $\alpha=\mathrm{roots}(p)$; and by
Lemma~\ref{lem:simple}(i) each $\alpha^\epsilon$ has \emph{distinct} entries, so
$\qfourInfo(T_\epsilon p)=\|\qfourScore(\alpha^\epsilon)\|^2$ is finite for every
$\epsilon>0$. We distinguish two cases.

\smallskip\noindent\emph{Case 1: $p$ has simple roots.} Then $\alpha$ has distinct
entries. On the open set $U=\{\xi\in\qfourR^n:\ \xi_i\ne\xi_j\ \forall i\ne j\}$ the
score map $\xi\mapsto\qfourScore(\xi)=\bigl(\sum_{j\ne i}(\xi_i-\xi_j)^{-1}\bigr)_i$
is continuous (each summand is continuous and the denominators are nonzero on $U$),
so $\xi\mapsto\qfourInfo(\xi)=\|\qfourScore(\xi)\|^2$ is continuous and finite on $U$.
Since $\alpha\in U$ and $\alpha^\epsilon\to\alpha$, we have $\alpha^\epsilon\in U$ for
all small $\epsilon$ and
\[
\qfourInfo(T_\epsilon p)=\|\qfourScore(\alpha^\epsilon)\|^2
\xrightarrow[\epsilon\to0]{}\|\qfourScore(\alpha)\|^2=\qfourInfo(p)\in(0,\infty).
\]
(The limit is strictly positive: ordering $\alpha_1<\dots<\alpha_n$, the top score
$\qfourScore(\alpha)_n=\sum_{j<n}1/(\alpha_n-\alpha_j)>0$ for $n\ge2$, so
$\qfourInfo(p)>0$; for $n=1$ both sides are the empty sum, but then there is no
simple-vs-repeated distinction and $\qfourInfo(p)=0$, handled by the convention.)
As $t\mapsto1/t$ is continuous at the finite positive value $\qfourInfo(p)$, we
conclude $1/\qfourInfo(T_\epsilon p)\to1/\qfourInfo(p)$.

\smallskip\noindent\emph{Case 2: $p$ has a root of multiplicity $m\ge2$.} Then
$\qfourInfo(p)=\infty$, i.e.\ $1/\qfourInfo(p)=0$, and we must show
$\qfourInfo(T_\epsilon p)\to+\infty$. We give the direct quantitative argument
(which is exactly the gap blow-up estimate of Lemma~\ref{lem:gapblowup} applied to
the present family, and which makes the blow-up explicit and rate-independent).

Let $\lambda_0$ be a root of $p$ of multiplicity $m\ge2$, and let
$G\subseteq\{1,\dots,n\}$, $|G|=m$, be the set of indices with $\alpha_i=\lambda_0$.
Define the cluster diameter
\[
d_\epsilon:=\max_{i,j\in G}\bigl|\alpha^\epsilon_i-\alpha^\epsilon_j\bigr|>0 .
\]
Because $\alpha^\epsilon_i\to\alpha_i=\lambda_0$ for every $i\in G$, all pairwise
differences within $G$ tend to $0$, hence
\begin{equation}\label{eq:deps}
d_\epsilon\xrightarrow[\epsilon\to0]{}0 .
\end{equation}
(We do \emph{not} need the precise rate of $d_\epsilon$; in fact for the Nuij
operator one has $d_\epsilon=\Theta(\epsilon)$, but only $d_\epsilon\to0$ is used.)

We now bound $\qfourInfo(T_\epsilon p)$ from below using the cluster. Let
$\eta^\epsilon:=(\alpha^\epsilon_i)_{i\in G}\in\qfourR^m$ be the sub-vector of
cluster coordinates; it has distinct entries (as $\alpha^\epsilon$ does) and
diameter $d_\epsilon$. Split each cluster score into internal and external parts:
for $i\in G$,
\[
\qfourScore(\alpha^\epsilon)_i
=\underbrace{\sum_{l\in G,\,l\ne i}\frac1{\alpha^\epsilon_i-\alpha^\epsilon_l}}_{I^\epsilon_i}
+\underbrace{\sum_{l\notin G}\frac1{\alpha^\epsilon_i-\alpha^\epsilon_l}}_{E^\epsilon_i}.
\]
For $l\notin G$ we have $\alpha^\epsilon_l\to\alpha_l\ne\lambda_0$ and
$\alpha^\epsilon_i\to\lambda_0$, so the external denominators stay bounded away from
$0$; hence there is a constant $C<\infty$ (independent of $\epsilon$, for small
$\epsilon$) with $|E^\epsilon_i|\le C$ for all $i\in G$. Moreover
$(I^\epsilon_i)_{i\in G}=\qfourScore(\eta^\epsilon)$ is exactly the $m$-point score
vector of $\eta^\epsilon$, so $\sum_{i\in G}(I^\epsilon_i)^2=\qfourInfo(\eta^\epsilon)$.
By translation invariance and the $1/d$ scaling of scores under dilation (each
difference of $\eta^\epsilon$ has absolute value $\le d_\epsilon$, and after
normalising $\eta^\epsilon=\mu_\epsilon\mathbf1+d_\epsilon\zeta^\epsilon$ with
$\zeta^\epsilon\in K_m'$ the unit-diameter configurations of
Lemma~\ref{lem:gapblowup}),
\[
\qfourInfo(\eta^\epsilon)=\frac{1}{d_\epsilon^2}\,\qfourInfo(\zeta^\epsilon)
\ \ge\ \frac{\rho_m}{d_\epsilon^2},\qquad \rho_m:=\inf_{\zeta\in K_m'}\qfourInfo(\zeta)\ge1>0,
\]
the bound $\rho_m\ge1$ being established in Lemma~\ref{lem:gapblowup} (the top
normalised score is a sum of $m-1\ge1$ positive terms each $\ge1$). Combining with
the elementary inequality $(a+b)^2\ge\tfrac12a^2-b^2$,
\[
\qfourInfo(T_\epsilon p)\ \ge\ \sum_{i\in G}\qfourScore(\alpha^\epsilon)_i^2
=\sum_{i\in G}(I^\epsilon_i+E^\epsilon_i)^2
\ \ge\ \tfrac12\sum_{i\in G}(I^\epsilon_i)^2-\sum_{i\in G}(E^\epsilon_i)^2
\ \ge\ \frac{\rho_m}{2\,d_\epsilon^2}-mC^2 .
\]
By \eqref{eq:deps}, $d_\epsilon\to0$, so the right-hand side tends to $+\infty$;
therefore $\qfourInfo(T_\epsilon p)\to+\infty$ and
$1/\qfourInfo(T_\epsilon p)\to0=1/\qfourInfo(p)$.

\smallskip
In both cases $1/\qfourInfo(T_\epsilon p)\to1/\qfourInfo(p)$, proving the lemma.
\end{proof}

\noindent\emph{Remark on the direction of semicontinuity.} In Case~2 it is the
\emph{lower} bound $\qfourInfo(T_\epsilon p)\ge\rho_m/(2d_\epsilon^2)-mC^2$ that is
needed and that we use: it forces $\qfourInfo(T_\epsilon p)\to+\infty$, equivalently
$1/\qfourInfo(T_\epsilon p)\to0$. Lower semicontinuity of $\qfourInfo$ is exactly
this lower bound and is therefore the correct (not the wrong) direction here, since
the target value $\qfourInfo(p)=+\infty$ can only be approached from below. No
upper bound on $\qfourInfo(T_\epsilon p)$ is required in Case~2, and in Case~1 the
full continuity established there supplies convergence directly.

\medskip
Applying Lemma~\ref{lem:Bconv} to $p$ and to $q$, and adding,
\[
b_\epsilon:=\frac{1}{\qfourInfo(T_\epsilon p)}+\frac{1}{\qfourInfo(T_\epsilon q)}
\;\xrightarrow[\epsilon\to0]{}\;
L:=\frac{1}{\qfourInfo(p)}+\frac{1}{\qfourInfo(q)} .
\]
(Here $L\in[0,\infty)$; the right-hand side of the theorem is well defined since
each reciprocal lies in $[0,\infty)$, with the convention $1/\infty=0$.)

\medskip
\noindent\textit{Step C: upper semicontinuity of the left-hand side.}
By Step~A, $\gamma^\epsilon\to\gamma$, and Lemma~\ref{lem:simple}(ii) (LSC applied
to $\gamma^\epsilon\to\gamma$) gives
\[
\liminf_{\epsilon\to 0}\qfourInfo\bigl(T_\epsilon^2(p\qfourbn q)\bigr)\;\ge\;\qfourInfo(p\qfourbn q).
\]
The map $t\mapsto 1/t$ is decreasing on $(0,\infty]$ (with $1/\infty=0$); applying
it to the inequality $\liminf \qfourInfo(\cdots)\ge\qfourInfo(p\qfourbn q)>0$ and
using that $\liminf$ of a quantity bounds its reciprocal's $\limsup$ from the
correct side gives
\[
a_\epsilon:=\frac{1}{\qfourInfo\bigl(T_\epsilon^2(p\qfourbn q)\bigr)},\qquad
\limsup_{\epsilon\to 0}a_\epsilon\;\le\;\frac{1}{\qfourInfo(p\qfourbn q)}.
\]
(Explicitly: if $\liminf_\epsilon\Phi_\epsilon\ge\Phi_0\in(0,\infty]$ then for any
$\Phi_0'<\Phi_0$ we have $\Phi_\epsilon\ge\Phi_0'$ for all small $\epsilon$, so
$a_\epsilon\le 1/\Phi_0'$ eventually, whence $\limsup a_\epsilon\le 1/\Phi_0'$;
letting $\Phi_0'\uparrow\Phi_0$ yields $\limsup a_\epsilon\le 1/\Phi_0$.)

\medskip
\noindent\textit{Assembling the chain.}
By \eqref{eq:eps-ineq}, $a_\epsilon\ge b_\epsilon$ for all $\epsilon>0$. By Step~B,
$b_\epsilon\to L$. For \emph{every} sequence $\epsilon_k\downarrow 0$ we then have
$a_{\epsilon_k}\ge b_{\epsilon_k}\to L$, so $\liminf_{k}a_{\epsilon_k}\ge L$;
since this holds for all such sequences, $\liminf_{\epsilon\to0}a_\epsilon\ge L$.
Combining with Step~C and $\liminf\le\limsup$,
\[
\frac{1}{\qfourInfo(p\qfourbn q)}
\;\ge\;\limsup_{\epsilon\to 0}a_\epsilon
\;\ge\;\liminf_{\epsilon\to 0}a_\epsilon
\;\ge\;L
\;=\;\frac{1}{\qfourInfo(p)}+\frac{1}{\qfourInfo(q)}.
\]
(The use of $1/\qfourInfo(p\qfourbn q)$ in the first inequality presumes
$\qfourInfo(p\qfourbn q)>0$. This is automatic for $n\ge 2$: if $p\qfourbn q$ has
a repeated root then $\qfourInfo(p\qfourbn q)=\infty$; if its roots
$\gamma_1\le\dots\le\gamma_n$ are distinct, then the largest root $\gamma_n$
satisfies $\qfourScore(\gamma)_n=\sum_{j<n}\frac{1}{\gamma_n-\gamma_j}>0$, since
every summand is positive, whence $\qfourInfo(p\qfourbn q)\ge\qfourScore(\gamma)_n^2>0$.
For $n=1$ both sides of the theorem are $+\infty$ under the convention $1/0=+\infty$,
so the statement holds trivially.)

This establishes the theorem for general monic real-rooted $p,q$. Thus we may assume henceforth that $p$, $q$,
and $p\qfourbn q$ all have simple roots. Then
$\alpha,\beta,\gamma=\qfourRootMap(\alpha,\beta)$ have distinct entries. We now
verify that, on this simple-root locus, $\qfourRootMap$ is in fact $C^\infty$, so
that its Jacobian $\qfourJac$ and the coordinate Hessians
$H^{(i)}:=\qfourHess(\qfourRootMap)_i$ used below exist classically.

\begin{qfourlemma}[smoothness of $\qfourRootMap$ on the simple-root locus]\label{lem:smooth}
Let $(\alpha_0,\beta_0)\in\qfourR^n\times\qfourR^n$ be such that $p\qfourbn q$
has $n$ distinct (real) roots $\gamma_1>\dots>\gamma_n$. Then there is a
neighbourhood $U$ of $(\alpha_0,\beta_0)$ in $\qfourR^n\times\qfourR^n$ and
$C^\infty$ functions $\gamma_1,\dots,\gamma_n:U\to\qfourR$ such that, for every
$(\alpha,\beta)\in U$, the numbers $\gamma_1(\alpha,\beta)>\dots>\gamma_n(\alpha,\beta)$
are exactly the ordered roots of $p\qfourbn q$; in particular
$\qfourRootMap=(\gamma_1,\dots,\gamma_n)$ is $C^\infty$ on the simple-root locus.
\end{qfourlemma}

\begin{proof}
By the explicit coefficient formula for $\qfourbn$ (equivalently \eqref{eq:walsh}),
the coefficients $c_0=1,c_1,\dots,c_n$ of
$(p\qfourbn q)(x)=\sum_{k=0}^n c_k x^{n-k}$ are polynomials in the elementary
symmetric functions of $\alpha$ and of $\beta$, hence \emph{polynomial}—in
particular $C^\infty$—functions of $(\alpha,\beta)$. Define
\[
F:\qfourR\times(\qfourR^n\times\qfourR^n)\to\qfourR,\qquad
F(x;\alpha,\beta):=(p\qfourbn q)(x)=\sum_{k=0}^n c_k(\alpha,\beta)\,x^{n-k},
\]
which is $C^\infty$ jointly in all its arguments. Fix an index $i$ and the root
$x_0:=\gamma_i$ of $F(\,\cdot\,;\alpha_0,\beta_0)$. Because the roots are
distinct, $x_0$ is a \emph{simple} root, so
\[
\partial_x F(x_0;\alpha_0,\beta_0)=(p\qfourbn q)'(x_0)
=\prod_{l\ne i}(x_0-\gamma_l)\neq 0 .
\]
The implicit function theorem applied to the equation $F(x;\alpha,\beta)=0$ at
$(x_0;\alpha_0,\beta_0)$ therefore yields a neighbourhood $U_i$ of
$(\alpha_0,\beta_0)$ and a unique $C^\infty$ function
$\gamma_i:U_i\to\qfourR$ with $\gamma_i(\alpha_0,\beta_0)=x_0$ and
$F(\gamma_i(\alpha,\beta);\alpha,\beta)\equiv 0$ on $U_i$. Doing this for each of
the $n$ distinct roots and intersecting the neighbourhoods gives a common
neighbourhood $U=\bigcap_i U_i$ on which $\gamma_1,\dots,\gamma_n$ are $C^\infty$
and pairwise distinct (after possibly shrinking $U$, by continuity of the
$\gamma_i$ and $\gamma_1(\alpha_0,\beta_0)>\dots>\gamma_n(\alpha_0,\beta_0)$).
These $n$ continuous functions take $n$ distinct values that are all roots of the
degree-$n$ polynomial $p\qfourbn q$, hence they are \emph{all} of its roots; in
particular they remain its ordered roots throughout $U$, so on $U$ the ordered
root map equals $(\gamma_1,\dots,\gamma_n)$ and is $C^\infty$.
\end{proof}

\noindent Henceforth, on the simple-root locus, $\qfourRootMap$ is $C^\infty$ by
Lemma~\ref{lem:smooth}; let $\qfourJac$ denote the Jacobian of $\qfourRootMap$ at
$(\alpha,\beta)$, and write $H^{(i)}:=\qfourHess(\qfourRootMap)_i$ for the
(symmetric) Hessian of the $i$-th coordinate function. Both exist classically and
are used freely below.

\section*{Step 1: score vectors and the Jacobian}

\begin{qfourobservation}\label{obs:scores}
For all $a,b\ge0$,
\[
\qfourJac\bigl(a\,\qfourScore(\alpha),\,b\,\qfourScore(\beta)\bigr)\;=\;(a+b)\,\qfourScore(\gamma).
\]
\end{qfourobservation}

\begin{proof}
Set $p_t=\exp(-\tfrac t2\partial_x^2)p$ with ordered roots $\alpha(t)$, and
define $q_t,\beta(t)$ and $r_t,\gamma(t)$ analogously.
By Corollary~\ref{cor:heat-commutes} (with $s=t$),
\[
r_{(a+b)t}\;=\;p_{at}\qfourbn q_{bt},
\]
so $\gamma((a+b)t)=\qfourRootMap(\alpha(at),\beta(bt))$. Differentiating at
$t=0$ and using the chain rule,
\[
(a+b)\,\gamma'(0)\;=\;\qfourJac\!\begin{pmatrix} a\,\alpha'(0)\\ b\,\beta'(0)\end{pmatrix}.
\]
Lemma~\ref{lem:heat} identifies $\alpha'(0)=\qfourScore(\alpha)$, etc., giving the
claim.
\end{proof}

\section*{Step 2: the Jacobian is a contraction on mean-zero vectors}

This is the substance of the proof.

\begin{qfourproposition}\label{prop:contraction}
If $u,v\in\qfourR^n$ are each orthogonal to $\mathbf 1_n$, then
\[
\|\qfourJac(u,v)\|^2\;\le\;\|u\|^2+\|v\|^2.
\]
\end{qfourproposition}

Write $\gamma_i=(\qfourRootMap)_i(\alpha,\beta)$ for the coordinate functions and
$H^{(i)}:=\qfourHess(\qfourRootMap)_i$. Let
$V=\{(u,v)\in\qfourR^n\times\qfourR^n:u^\ast\mathbf 1_n=v^\ast\mathbf 1_n=0\}$.

\begin{qfourlemma}[Hessian identity]\label{lem:hess}
For all $w\in V$,
\[
w^\ast \qfourJac^\ast\qfourJac\, w
\;=\;w^\ast\Bigl(I_n\oplus I_n-\sum_{i=1}^n\gamma_i H^{(i)}\Bigr)w.
\]
\end{qfourlemma}

\begin{proof}
Fix $w=(u,v)\in V$ and set $\alpha(t)=\alpha+tu$, $\beta(t)=\beta+tv$,
$\gamma(t)=\qfourRootMap(\alpha(t),\beta(t))$. By the additivity of variance
(Proposition~\ref{prop:props}(i)),
\[
m_2(\gamma(t))-m_1(\gamma(t))^2
=m_2(\alpha(t))+m_2(\beta(t))-m_1(\alpha(t))^2-m_1(\beta(t))^2 .
\]
Since $(u,v)\in V$, the means $m_1(\alpha(t)),m_1(\beta(t))$ — hence
$m_1(\gamma(t))$ — are constant in $t$. Differentiating twice at $t=0$,
\[
\partial_t^2 m_2(\gamma(t))\big|_0=\partial_t^2\bigl(m_2(\alpha(t))+m_2(\beta(t))\bigr)\big|_0 .
\]
To evaluate the left side, we expand $\gamma_i(t)$ in a Taylor series at $t=0$.
Since $\alpha(t)=\alpha+tu$ and $\beta(t)=\beta+tv$ are linear in $t$, the smooth
function $\gamma_i(t)=(\qfourRootMap)_i(\alpha(t),\beta(t))$ expands as
\[
\gamma_i(t) = \gamma_i + t\,(\qfourJac w)_i + \tfrac{t^2}{2}\,w^\ast H^{(i)}w + O(t^3),
\]
where $(\qfourJac w)_i = \partial_t\gamma_i\big|_0$ is the $i$-th component of the
Jacobian image of $w$, and $w^\ast H^{(i)}w = \partial_t^2\gamma_i\big|_0$ is the
Hessian quadratic form. Squaring and retaining terms through order $t^2$:
\[
\gamma_i(t)^2 = \gamma_i^2 + 2t\,\gamma_i(\qfourJac w)_i
  + t^2\bigl[(\qfourJac w)_i^2 + \gamma_i\,w^\ast H^{(i)}w\bigr] + O(t^3).
\]
Taking $\partial_t^2\big|_{t=0}$ extracts twice the $t^2$-coefficient:
\[
\partial_t^2\bigl(\gamma_i(t)^2\bigr)\big|_0
  = 2\bigl[(\qfourJac w)_i^2 + \gamma_i\,w^\ast H^{(i)}w\bigr].
\]
Summing over $i$, dividing by $n$, and using $\sum_i(\qfourJac w)_i^2
=\|\qfourJac w\|^2=w^\ast\qfourJac^\ast\qfourJac\,w$:
\[
\partial_t^2 m_2(\gamma(t))\big|_0
  = \tfrac{1}{n}\sum_{i=1}^n \partial_t^2\bigl(\gamma_i(t)^2\bigr)\big|_0
  = \tfrac{2}{n}\Bigl(w^\ast\qfourJac^\ast\qfourJac\,w
    + \sum_{i=1}^n\gamma_i\,w^\ast H^{(i)}w\Bigr).
\]
For the right side, since $\alpha(t)=\alpha+tu$ is linear,
$m_2(\alpha(t))=\tfrac1n\sum_i(\alpha_i+tu_i)^2
=m_2(\alpha)+\tfrac{2t}{n}\alpha^\ast u+\tfrac{t^2}{n}\|u\|^2$,
so $\partial_t^2 m_2(\alpha(t))\big|_0=\tfrac{2}{n}\|u\|^2$;
likewise $\partial_t^2 m_2(\beta(t))\big|_0=\tfrac{2}{n}\|v\|^2$.
Hence the right side equals $\tfrac{2}{n}(\|u\|^2+\|v\|^2)=\tfrac{2}{n}w^\ast w$.
Equating left and right sides and cancelling $\tfrac{2}{n}$:
\[
w^\ast\qfourJac^\ast\qfourJac\,w + \sum_{i=1}^n\gamma_i\,w^\ast H^{(i)}w = w^\ast w,
\]
which rearranges to the claimed identity.
\end{proof}

The minus sign means we must show $\sum_i\gamma_i H^{(i)}\qfourPSD$ on $V$. This is
where real-rootedness of $\qfourbn$ is used in an essential way, via the theory of
hyperbolic polynomials. We first recall the relevant definitions, then state the
convexity theorem we use (with a proof of the part we need), and finally verify
its hypotheses for the Walsh polynomial \eqref{eq:walsh}.

\subsection*{Hyperbolic polynomials and their eigenvalues}

\begin{qfourdefinition}[hyperbolicity]\label{def:hyp}
Let $f\in\qfourR[z_1,\dots,z_m]$ be homogeneous of degree $d$ and let
$e\in\qfourR^m$ be a vector with $f(e)\ne0$. We say $f$ is \emph{hyperbolic in
the direction $e$} if for every $z\in\qfourR^m$ the univariate polynomial
\[
t\;\longmapsto\;f(z-te)
\]
has only real roots. (By homogeneity its degree in $t$ equals $d$ and its
leading coefficient is $(-1)^d f(e)\ne0$, so it has exactly $d$ roots counted
with multiplicity.) For $z\in\qfourR^m$ we list these $d$ real roots, in
decreasing order, as
\[
\lambda_1(z)\;\ge\;\lambda_2(z)\;\ge\;\dots\;\ge\;\lambda_d(z),
\]
and call them the \emph{hyperbolic eigenvalues of $z$} (relative to $e$). Thus
\[
f(z-te)\;=\;f(e)\prod_{i=1}^d\bigl(\lambda_i(z)-t\bigr),\qquad
f(z)\;=\;f(e)\prod_{i=1}^d\lambda_i(z).
\]
\end{qfourdefinition}

\noindent\emph{Motivating example.} For $f(z)=\det Z$ on the space of
$d\times d$ real symmetric matrices $Z$ (a homogeneous polynomial of degree $d$
in the $\binom{d+1}{2}$ entries) and $e=I$, hyperbolicity in $e$ is the
statement that $\det(Z-tI)$ has only real roots for symmetric $Z$ -- the
spectral theorem -- and $\lambda_1(Z)\ge\dots\ge\lambda_d(Z)$ are exactly the
eigenvalues of $Z$. This is the prototype the terminology imitates; the
classical fact ``$Z\mapsto\sum_{i\le k}\lambda_i(Z)$ is convex'' (Ky Fan) is the
special case of Theorem~\ref{thm:bgls} below.

\smallskip
We will use one structural fact, due to G\aa rding, in the proof of the
convexity theorem.

\begin{qfourlemma}[G\aa rding's cone; \cite{Garding}]\label{lem:garding}
Let $f$ be hyperbolic of degree $d$ in direction $e$, with eigenvalues
$\lambda_1\ge\dots\ge\lambda_d$ as in Definition~\ref{def:hyp}. Then:
\begin{enumerate}\setlength\itemsep{2pt}
\item the smallest eigenvalue
$\lambda_d(z)=\min_i\lambda_i(z)$ is a concave function of $z$, and the largest
eigenvalue $\lambda_1(z)=\max_i\lambda_i(z)$ is a convex function of $z$;
\item more generally $f$ is hyperbolic in direction $e'$ for every $e'$ in the
open ``hyperbolicity cone'' $\Lambda=\{z:\lambda_d(z)>0\}$, which is an open
convex cone containing $e$.
\end{enumerate}
\end{qfourlemma}

\noindent (Part (i) is all we strictly need; it is classical and is, e.g.,
\cite[\S2]{BGLS} or \cite{Garding}. We include the short proof of the convex
combination statement we actually invoke directly in the proof of
Theorem~\ref{thm:bgls}, so the argument is self-contained up to G\aa rding's
theorem.)

\subsection*{The convexity theorem of Bauschke--G\"uler--Lewis--Sendov}

\begin{qfourtheorem}[Bauschke--G\"uler--Lewis--Sendov \cite{BGLS}]\label{thm:bgls}
Let $f\in\qfourR[z_1,\dots,z_m]$ be homogeneous of degree $d$ and hyperbolic in
the direction $e$, with hyperbolic eigenvalues
$\lambda_1(z)\ge\dots\ge\lambda_d(z)$ as in Definition~\ref{def:hyp}. Then for
each $k\in\{1,\dots,d\}$ the \emph{partial eigenvalue sum}
\[
\sigma_k(z)\;:=\;\sum_{i=1}^k\lambda_i(z)
\]
is a convex function of $z\in\qfourR^m$. (In particular, if $z\mapsto\sigma_k(z)$
is twice differentiable at a point $z_0$, then $\qfourHess\sigma_k(z_0)\qfourPSD$.)
\end{qfourtheorem}

\begin{proof}[Proof sketch of the convexity, following \cite{BGLS}]
We indicate why $\sigma_k$ is convex; the argument is that of \cite{BGLS}, and we
record the two ingredients: G\aa rding's theorem (Step~1) for the $k=1$ base case,
and BGLS's variational sup-representation (Step~2) for general $k$.

\smallskip\noindent\emph{Step 1: $\sigma_1=\lambda_1$ is convex.} This is
Lemma~\ref{lem:garding}(i): the largest eigenvalue of a hyperbolic polynomial is
a convex function of $z$ (it is the prototypical statement of G\aa rding's
theory; the convexity of the hyperbolicity cone $\Lambda=\{\lambda_d>0\}$ is
equivalent to the concavity of $\lambda_d$, and applying this to $-z$ gives
convexity of $\lambda_1$).

\smallskip\noindent\emph{Step 2: $\sigma_k$ is a pointwise supremum of affine
functions of $z$, hence convex.}
Bauschke--G\"uler--Lewis--Sendov establish \cite[\S4--5]{BGLS} that for each
$k\in\{1,\dots,d\}$ the partial eigenvalue sum $\sigma_k(z)=\sum_{i\le k}\lambda_i(z)$
can be written as a pointwise supremum over a family of affine (linear) functions of
$z$. The prototype is Ky~Fan's variational identity for $d\times d$ symmetric matrices
$Z$:
\[
\sigma_k(Z)=\max_{\substack{P\in\mathbb R^{d\times d}\\P^2=P,\;P^\top=P,\;\mathrm{tr}(P)=k}}
\mathrm{tr}(PZ),
\]
where each function $Z\mapsto\mathrm{tr}(PZ)$ is linear in $Z$; BGLS generalize this
variational principle to eigenvalues of an arbitrary hyperbolic polynomial $f$.
Since a pointwise supremum of affine functions is automatically convex (each affine
function is convex and convexity is preserved under pointwise sup), $\sigma_k$ is
convex.  (For $k=d$: $\sigma_d(z)=\sum_i\lambda_i(z)$ equals the coefficient of
$t^{d-1}$ in $f(z-te)/f(e)$ up to sign and is a \emph{linear} function of $z$ by
homogeneity, in particular convex — matching the affine behaviour of $\sigma_n$ used
in Corollary~\ref{cor:psd}.)

\smallskip
A convex function that is finite on all of $\qfourR^m$ has a positive
semidefinite Hessian at every point of twice-differentiability (and PSD Hessian
in the distributional sense everywhere); this is the only consequence we use.
For a complete treatment of Steps~1--2, including the variational principle for
arbitrary hyperbolic polynomials in Step~2, see \cite[Thms.~1--2 and \S4--5]{BGLS}.
\end{proof}

\noindent\emph{Remark on what we actually use.} The only consequence of
Theorem~\ref{thm:bgls} invoked below is: \emph{each $\sigma_k$ is a convex
function of $z$, hence has positive semidefinite Hessian wherever it is twice
differentiable.} At a point $z_0$ where the top $k$ eigenvalues
$\lambda_1(z_0)>\lambda_{k+1}(z_0)$ are strictly separated (a simple-spectrum /
spectral-gap condition), $\sigma_k$ is real-analytic near $z_0$ by the implicit
function theorem applied to the relevant factor of $f(z-te)$, so the Hessian
exists; this is the situation we are in (the roots $\gamma_i$ are distinct after
the reduction to simple roots). Where the spectrum is degenerate, convexity
still gives $\qfourHess\sigma_k\qfourPSD$ in the distributional sense, which is
all that is needed for the averaging in Corollary~\ref{cor:psd}.

\subsection*{The Walsh polynomial is hyperbolic in $e_1$}

\begin{qfourcorollary}\label{cor:psd}
For any reals $d_1\ge\dots\ge d_n$, the matrix $\sum_{i=1}^n d_i H^{(i)}$ is PSD.
In particular, taking $d_i=\gamma_i$ (the decreasingly ordered roots of
$p\qfourbn q$), the matrix $\sum_{i=1}^n\gamma_i H^{(i)}$ is PSD.
\end{qfourcorollary}

\begin{proof}
Work in $\qfourR^m$ with $m=2n+1$ and coordinates $z=(x,a_1,\dots,a_n,b_1,\dots,b_n)$,
and consider the \emph{Walsh polynomial}
\[
f(z)\;=\;f(x,a,b)\;=\;\frac{1}{n!}\sum_{\pi\in S_n}\prod_{i=1}^n\bigl(x-a_i-b_{\pi(i)}\bigr).
\]

\smallskip\noindent\textit{$f$ is homogeneous of degree $n$.} Each factor
$x-a_i-b_{\pi(i)}$ is a homogeneous linear form in $z$, so each product is
homogeneous of degree $n$ and so is their average.

\smallskip\noindent\textit{$f$ is hyperbolic in $e_1=(1,0,\dots,0)$.} First,
$f(e_1)=\frac1{n!}\sum_\pi\prod_i(1-0-0)=1\ne0$. Next we verify the root
condition of Definition~\ref{def:hyp}: for every real $z=(x,a,b)$ the univariate
polynomial $t\mapsto f(z-te_1)=f(x-t,a,b)$ has only real roots. Indeed, by
\eqref{eq:walsh},
\[
x'\;\longmapsto\;f(x',a,b)\;=\;\Bigl(\textstyle\prod_i(x'-a_i)\Bigr)\;\qfourbn\;
\Bigl(\textstyle\prod_j(x'-b_j)\Bigr),
\]
which is the finite free convolution of two real-rooted polynomials and hence is
itself real-rooted by Walsh's reality theorem (Lemma~\ref{lem:walsh}); write its
roots as $\rho_1\ge\dots\ge\rho_n\in\qfourR$. Since $f(z-te_1)=f(x-t,a,b)$, its
roots in $t$ are exactly $t=x-\rho_i$, all real. Thus $f$ is hyperbolic in
$e_1$, with hyperbolic eigenvalues
\[
\lambda_i(z)\;=\;x-\rho_{n+1-i}(a,b)\qquad(\text{ordered }\lambda_1\ge\dots\ge\lambda_n).
\]
Both homogeneity and hyperbolicity are now verified, so
Theorem~\ref{thm:bgls} applies: each
\[
\sigma_k(z)\;=\;\sum_{i=1}^k\lambda_i(z)
\]
is a convex function of $z$, with $\qfourHess\sigma_k(z)\qfourPSD$ wherever twice
differentiable.

\smallskip\noindent\textit{Abel summation.} Fix reals $c_1\ge\dots\ge c_n$ and
set
\[
L(z)\;:=\;\sum_{i=1}^n c_i\,\lambda_i(z).
\]
Because $c_1\ge\dots\ge c_n$, the Abel (summation-by-parts) identity gives
\[
L(z)\;=\;\sum_{i=1}^n c_i\lambda_i(z)
\;=\;\sum_{k=1}^{n-1}(c_k-c_{k+1})\,\sigma_k(z)\;+\;c_n\,\sigma_n(z),
\]
which one checks by collecting, on the right-hand side, the coefficient of each
$\lambda_i$ (recall $\sigma_k=\sum_{j\le k}\lambda_j$, so $\lambda_i$ appears in
$\sigma_k$ exactly when $k\ge i$): for $i<n$ the coefficient is
$\sum_{k=i}^{n-1}(c_k-c_{k+1})+c_n=c_i-c_n+c_n=c_i$, and for $i=n$ it is $c_n$.
The weights $c_k-c_{k+1}\ge0$ for $k=1,\dots,n-1$ since $c_1\ge\dots\ge c_n$.

\smallskip\noindent\textit{The top term is affine.} The full sum of eigenvalues
is
\[
\sigma_n(z)\;=\;\sum_{i=1}^n\lambda_i(z)\;=\;\sum_{i=1}^n(x-\rho_i(a,b))
\;=\;nx-\sum_{i=1}^n\rho_i(a,b)
\;=\;nx-\Bigl(\sum_i a_i+\sum_j b_j\Bigr),
\]
where the last equality is the additivity of means
(Proposition~\ref{prop:props}(i)): the sum of the roots of $p\qfourbn q$ equals
$\sum_i\alpha_i+\sum_j\beta_j$. Hence $\sigma_n$ is \emph{affine} in $z$, so
$\qfourHess\sigma_n\equiv0$.

\smallskip\noindent\textit{Conclusion (Hessian of $L$ is PSD).} Combining,
\[
\qfourHess L(z)
\;=\;\sum_{k=1}^{n-1}(c_k-c_{k+1})\,\qfourHess\sigma_k(z)\;+\;c_n\cdot0
\;=\;\sum_{k=1}^{n-1}(c_k-c_{k+1})\,\qfourHess\sigma_k(z),
\]
a nonnegative combination of PSD matrices (Theorem~\ref{thm:bgls}, since
$c_k-c_{k+1}\ge0$), hence $\qfourHess L(z)\qfourPSD$ for every $z$. Crucially, the
only possibly-negative coefficient $c_n$ multiplies the \emph{affine} term
$\sigma_n$, whose Hessian vanishes, so no sign hypothesis on $c_n$ is required.

\smallskip\noindent\textit{Identifying $\qfourHess L$ with $\sum_i c_i' H^{(i)}$ via restriction to $\{x=0\}$.}
We now relate $\qfourHess L$, computed in the full space $\qfourR^m$ with
$m=2n+1$, to the Hessians $H^{(i)}=\qfourHess(\qfourRootMap)_i$, which live on the
$2n$-dimensional $(a,b)$-space. The two objects have different sizes
($(2n{+}1)\times(2n{+}1)$ versus $2n\times2n$), so the passage between them is a
genuine \emph{restriction} step, which we now carry out explicitly.

Recall from the hyperbolicity computation that, with $\rho_1\ge\dots\ge\rho_n$
denoting the roots of the univariate polynomial
\[
x'\;\longmapsto\;f(x',a,b)\;=\;\Bigl(\textstyle\prod_i(x'-a_i)\Bigr)\;\qfourbn\;
\Bigl(\textstyle\prod_j(x'-b_j)\Bigr),
\]
the $i$-th (decreasingly ordered) hyperbolic eigenvalue of $f$ is
\begin{equation}\label{eq:lambda-eq-rho}
\lambda_i(z)\;=\;\lambda_i(x,a,b)\;=\;x-\rho_{\,n+1-i}(a,b)\qquad(1\le i\le n).
\end{equation}
This is the explicit identification requested in the gap: the BGLS eigenvalue
$\lambda_i$, viewed as a function on the full $(2n{+}1)$-dimensional space, is
the root $\rho_{n+1-i}(a,b)$ of $p\qfourbn q$ \emph{shifted by $x$}. (The sign
and the index reversal $i\mapsto n+1-i$ come solely from the conventions of
Definition~\ref{def:hyp}, where the eigenvalues are the roots of
$t\mapsto f(z-te_1)=f(x-t,a,b)$, i.e.\ $t=x-\rho_j$, listed decreasingly.)

\smallskip\noindent\emph{Per-eigenvalue Hessian blocks.}
Since $\rho_1\ge\dots\ge\rho_n$ are precisely the decreasingly ordered roots of
$p\qfourbn q$ and $H^{(\cdot)}$ is the Hessian (in the root parameters $(a,b)$)
of the corresponding root coordinate function, we have
\[
\qfourHess_{(a,b)}\rho_j(a,b)\;=\;H^{(j)}\qquad(1\le j\le n),
\]
where $H^{(j)}$ is read off in the \emph{decreasing} labelling. (Two remarks on
labelling are in order. First, this is consistent with the definition
$H^{(i)}:=\qfourHess(\qfourRootMap)_i$ used in Lemma~\ref{lem:hess}: there
$(\qfourRootMap)_i$ is the $i$-th ordered root of $p\qfourbn q$, and reversing
the ordering convention only relabels the family $\{H^{(j)}\}_{j=1}^n$ by the
permutation $j\mapsto n+1-j$. Second, this relabelling is immaterial for the
quantity we ultimately feed into Lemma~\ref{lem:hess}, namely
$\sum_i\gamma_i H^{(i)}$ with $\gamma_i=(\qfourRootMap)_i$: in that sum the
weight $\gamma_i$ and the matrix $H^{(i)}$ carry the \emph{same} index, so the
sum is invariant under any joint relabelling of the roots and is in particular
the same whether the roots are ordered increasingly or decreasingly. We will
therefore prove the PSD statement in the decreasing convention and transfer it
to the sum $\sum_i\gamma_iH^{(i)}$ at the end.)

By \eqref{eq:lambda-eq-rho}, each $\lambda_i$ is the sum of the linear function
$z\mapsto x$ and the function $(a,b)\mapsto-\rho_{n+1-i}(a,b)$ that does not
depend on $x$. Consequently its full Hessian on $\qfourR^{2n+1}$, in coordinates
$(x;a,b)$, is block-diagonal with vanishing $x$-block:
\begin{equation}\label{eq:Hesslambda-block}
\qfourHess_{(x,a,b)}\lambda_i
=\begin{pmatrix}0 & 0\\[2pt] 0 & \qfourHess_{(a,b)}\bigl[-\rho_{n+1-i}(a,b)\bigr]\end{pmatrix}
=\begin{pmatrix}0 & 0\\[2pt] 0 & -H^{(n+1-i)}\end{pmatrix}.
\end{equation}
In particular $\partial_x^2\lambda_i=\partial_x\partial_{a_j}\lambda_i
=\partial_x\partial_{b_j}\lambda_i=0$, and the lower-right $2n\times2n$ block of
$\qfourHess_{(x,a,b)}\lambda_i$ is exactly $-H^{(n+1-i)}$. Summing
\eqref{eq:Hesslambda-block} against the weights $c_i$, the same block structure
holds for $L=\sum_i c_i\lambda_i$; equivalently, $L$ is an \emph{affine} function
of $x$ (with constant $x$-coefficient $\sum_i c_i$) plus a function of $(a,b)$
alone.

\smallskip\noindent\emph{Block structure of the full Hessian.} Collecting
\eqref{eq:Hesslambda-block} for $L=\sum_i c_i\lambda_i$, and writing the
coordinates in the order $(x;a,b)$ with the orthogonal splitting
$\qfourR^{2n+1}=\qfourR e_1\oplus W$, $W=\{x=0\}\cong\qfourR^{2n}$ the
$(a,b)$-subspace, the full Hessian of $L$ is block-diagonal with a vanishing
$x$-block:
\begin{equation}\label{eq:HessL-block}
\qfourHess_{(x,a,b)}L(z)
\;=\;
\begin{pmatrix} 0 & 0\\[2pt] 0 & B(a,b)\end{pmatrix},
\qquad
B(a,b):=\qfourHess_{(a,b)}\bigl[L(0,a,b)\bigr].
\end{equation}
Here $B(a,b)$ is exactly the Hessian, in the $2n$ variables $(a,b)$, of the
\emph{restricted} function $\tilde L(a,b):=L(0,a,b)$ obtained by fixing $x=0$;
equivalently it is the lower-right $2n\times2n$ principal block of
$\qfourHess_{(x,a,b)}L$. (That $B$ does not depend on $x$ is again because $L$ is
affine in $x$.)

\smallskip\noindent\emph{The restricted Hessian is PSD.} The function
$x\mapsto L$ being convex on all of $\qfourR^{2n+1}$ (established above), its
restriction $\tilde L(a,b)=L(0,a,b)$ to the affine subspace $W=\{x=0\}$ is a
convex function on $\qfourR^{2n}$; consequently $B=\qfourHess\tilde L\qfourPSD$.
Equivalently, and purely matrix-theoretically: a principal submatrix of a PSD
matrix is PSD, and by \eqref{eq:HessL-block} the $(a,b)$-block $B$ is the
lower-right principal submatrix of $\qfourHess_{(x,a,b)}L\qfourPSD$; hence
$B\qfourPSD$. (Concretely, for $w\in\qfourR^{2n}$, embedding it as
$(0,w)\in\qfourR^{2n+1}$ gives $w^\ast B w=(0,w)^\ast\qfourHess L\,(0,w)\ge0$.)

\smallskip\noindent\emph{Computing the restricted Hessian.} Since each
$\lambda_i$ is affine in $x$,
\[
\qfourHess_{(a,b)}\bigl[\lambda_i(0,a,b)\bigr]
\;=\;\qfourHess_{(a,b)}\bigl[-\rho_{\,n+1-i}(a,b)\bigr]
\;=\;-\,H^{(n+1-i)},
\]
the constant ``$x$''-contribution dropping out under the $(a,b)$-Hessian. Summing
with weights $c_i$,
\[
B
=\qfourHess_{(a,b)}\tilde L
=\sum_{i=1}^n c_i\,\qfourHess_{(a,b)}\bigl[\lambda_i(0,a,b)\bigr]
=-\sum_{i=1}^n c_i\,H^{(n+1-i)}
=-\sum_{j=1}^n c_{\,n+1-j}\,H^{(j)},
\]
the last step being the reindexing $j=n+1-i$. Writing $c_j':=-c_{\,n+1-j}$,
the two displays above combine to give
\begin{equation}\label{eq:HessL-is-cprimeH}
\sum_{j=1}^n c_j'\,H^{(j)}\;=\;B\;=\;\qfourHess_{(a,b)}\tilde L\;\qfourPSD
\qquad\text{whenever }c_1\ge\dots\ge c_n .
\end{equation}
Reversing and negating the index, \eqref{eq:HessL-is-cprimeH} says precisely:
\[
\boxed{\ \text{for every }d_1\ge\dots\ge d_n,\qquad \sum_{j=1}^n d_j\,H^{(j)}\qfourPSD.\ }
\]
Indeed, given any $d_1\ge\dots\ge d_n$, set $c_i:=-d_{\,n+1-i}$; then
$c_1\ge\dots\ge c_n$ (because $d$ is decreasing), and $c_j'=-c_{n+1-j}=d_j$, so
\eqref{eq:HessL-is-cprimeH} yields $\sum_j d_j H^{(j)}\qfourPSD$. This is exactly the assertion of the Corollary.

\smallskip\noindent\textit{Evaluation at the relevant point.}
Finally, applying the boxed statement with $d_j=\gamma_j$ — the ordered roots
$\gamma_1\ge\dots\ge\gamma_n$ with $\gamma_j=\rho_j(\alpha,\beta)$ of $p\qfourbn q$ at the point
$(a,b)=(\alpha,\beta)$, where $H^{(j)}=\qfourHess\rho_j(\alpha,\beta)$ — gives
\[
\sum_{j=1}^n \gamma_j\,H^{(j)}\;\qfourPSD,
\]
which is the inequality used in Proposition~\ref{prop:contraction}. (When the
$\gamma_j$ are merely real but possibly coincident, the same conclusion holds
because $\qfourHess L\qfourPSD$ holds for the convex function $L$ in the
distributional sense; in any case the reduction to simple roots in
Lemma~\ref{lem:simple} lets us assume the $\gamma_j$ distinct, so the Hessians
$H^{(j)}$ exist classically.)
\end{proof}

\begin{proof}[Proof of Proposition~\ref{prop:contraction}]
Let $(u,v)\in V$. By Lemma~\ref{lem:hess},
\[
\|\qfourJac(u,v)\|^2=(u,v)^\ast \qfourJac^\ast\qfourJac\,(u,v)
=\|u\|^2+\|v\|^2-\sum_{i=1}^n\gamma_i\,(u,v)^\ast H^{(i)}(u,v).
\]
The roots $\gamma_1\ge\dots\ge\gamma_n$ are ordered, so
Corollary~\ref{cor:psd} (with $d_i=\gamma_i$) gives $\sum_i\gamma_i H^{(i)}\qfourPSD$,
whence $\sum_i\gamma_i(u,v)^\ast H^{(i)}(u,v)\ge0$. The claim follows.
\end{proof}

\section*{Step 3: Blachman's argument}

Before optimizing the Blachman inequality we record that the finite free Fisher
information is strictly positive on simple-root polynomials of degree $\ge2$; this
makes the optimal weights well defined and removes the degenerate case
$\qfourInfo=0$ flagged in the gap analysis.

\begin{qfourlemma}[strict positivity of $\qfourInfo$ on simple roots]\label{lem:positive}
Let $p$ be a monic polynomial of degree $n\ge2$ with simple roots
$\alpha_1,\dots,\alpha_n$. Then $\qfourScore(\alpha)\ne0$; equivalently
$\qfourInfo(p)=\|\qfourScore(\alpha)\|^2>0$.
\end{qfourlemma}

\begin{proof}
As computed in the proof of Lemma~\ref{lem:heat}, for each root $\alpha_i$ one has
$p'(\alpha_i)\ne0$ (simplicity) and
\[
\qfourScore(\alpha)_i\;=\;\sum_{j\ne i}\frac{1}{\alpha_i-\alpha_j}
\;=\;\frac{1}{2}\,\frac{p''(\alpha_i)}{p'(\alpha_i)} .
\]
Hence $\qfourScore(\alpha)=0$ would force $p''(\alpha_i)=0$ for all $n$ distinct
roots $\alpha_1,\dots,\alpha_n$. But $p$ is monic of degree $n$, so $p''$ is a
polynomial of degree $n-2$ with leading coefficient $n(n-1)\ne0$; in particular
$p''\not\equiv0$ and $p''$ has at most $n-2$ roots. Since $n-2<n$, $p''$ cannot
vanish at $n$ distinct points. This contradiction shows $\qfourScore(\alpha)\ne0$,
so $\qfourInfo(p)>0$.
\end{proof}

\begin{proof}[Proof of the Theorem (assuming simple roots)]
We prove the inequality for $p,q$ with simple roots; Lemma~\ref{lem:simple} then
extends to all real-rooted polynomials by the limiting argument in that lemma.

\smallskip\noindent\textit{The case $n=1$.} If $n=1$ the score vector is the empty
sum, so $\qfourScore(\alpha)=\qfourScore(\beta)=0$ and
$\qfourInfo(p)=\qfourInfo(q)=\qfourInfo(p\qfourbn q)=0$. Both sides of the claimed
inequality equal $1/0=+\infty$, so it holds (as $+\infty\ge+\infty$). We henceforth
assume $n\ge2$.

\smallskip\noindent\textit{The Blachman inequality.}
First, $\qfourScore(\alpha)\perp\mathbf 1_n$: indeed
$\sum_i\sum_{j\ne i}(\alpha_i-\alpha_j)^{-1}=0$ since each unordered pair
contributes $+$ and $-$ once. Likewise $\qfourScore(\beta)\perp\mathbf 1_n$. So
Proposition~\ref{prop:contraction} applies to
$(a\,\qfourScore(\alpha),\,b\,\qfourScore(\beta))$:
\[
\bigl\|\qfourJac\bigl(a\,\qfourScore(\alpha),b\,\qfourScore(\beta)\bigr)\bigr\|^2
\;\le\;a^2\|\qfourScore(\alpha)\|^2+b^2\|\qfourScore(\beta)\|^2 .
\]
Combining with Observation~\ref{obs:scores},
\begin{equation}\label{eq:blachman}
(a+b)^2\,\qfourInfo(p\qfourbn q)\;\le\;a^2\,\qfourInfo(p)+b^2\,\qfourInfo(q),
\end{equation}
valid for all $a,b\ge0$.

\smallskip\noindent\textit{Optimization (well defined because $\qfourInfo>0$).}
Since $p$ and $q$ have simple roots and $n\ge2$, Lemma~\ref{lem:positive} gives
$\qfourInfo(p)>0$ and $\qfourInfo(q)>0$, and both are finite (simple roots), so the
weights
\[
a:=\frac{1}{\qfourInfo(p)}\in(0,\infty),\qquad b:=\frac{1}{\qfourInfo(q)}\in(0,\infty)
\]
are well-defined finite positive reals. With this choice the right side
of~\eqref{eq:blachman} becomes
$a^2\qfourInfo(p)+b^2\qfourInfo(q)=\tfrac{1}{\qfourInfo(p)}+\tfrac{1}{\qfourInfo(q)}$,
while $(a+b)^2=\bigl(\tfrac{1}{\qfourInfo(p)}+\tfrac{1}{\qfourInfo(q)}\bigr)^2$, so
\[
\Bigl(\tfrac{1}{\qfourInfo(p)}+\tfrac{1}{\qfourInfo(q)}\Bigr)^2\qfourInfo(p\qfourbn q)
\;\le\;\tfrac{1}{\qfourInfo(p)}+\tfrac{1}{\qfourInfo(q)} .
\]
The factor $\tfrac{1}{\qfourInfo(p)}+\tfrac{1}{\qfourInfo(q)}$ is strictly positive,
so dividing through by it yields
\[
\Bigl(\tfrac{1}{\qfourInfo(p)}+\tfrac{1}{\qfourInfo(q)}\Bigr)\qfourInfo(p\qfourbn q)\le1,
\]
i.e.\ $\qfourInfo(p\qfourbn q)\le\bigl(\tfrac{1}{\qfourInfo(p)}+\tfrac{1}{\qfourInfo(q)}\bigr)^{-1}$.
If $\qfourInfo(p\qfourbn q)=0$ the desired inequality
$\tfrac{1}{\qfourInfo(p\qfourbn q)}=+\infty\ge\tfrac{1}{\qfourInfo(p)}+\tfrac{1}{\qfourInfo(q)}$
is trivial; otherwise $\qfourInfo(p\qfourbn q)>0$ and dividing the last display by it gives
\[
\frac{1}{\qfourInfo(p\qfourbn q)}\;\ge\;\frac{1}{\qfourInfo(p)}+\frac{1}{\qfourInfo(q)} .
\]
This completes the simple-root case.
\end{proof}

\begin{qfourobservation}[the degenerate case $\qfourInfo=0$ in general]\label{obs:degenerate}
For completeness we record how the Blachman inequality~\eqref{eq:blachman} alone
handles a vanishing score, without invoking Lemma~\ref{lem:positive}. Suppose
$\qfourInfo(p)=0$, i.e.\ $\qfourScore(\alpha)=0$ (by Lemma~\ref{lem:positive} this
does not occur for simple roots of degree $\ge2$, but the argument is robust and is
the one needed should one wish to bypass that lemma). Then for every $a\ge0$ we have
$a\,\qfourScore(\alpha)=0$, and~\eqref{eq:blachman} reads
\[
(a+b)^2\,\qfourInfo(p\qfourbn q)\;\le\;b^2\,\qfourInfo(q)\qquad\text{for all }a,b\ge0 .
\]
Fixing $b=1$ and letting $a\to+\infty$ forces $\qfourInfo(p\qfourbn q)=0$, hence
$\tfrac{1}{\qfourInfo(p\qfourbn q)}=+\infty\ge\tfrac{1}{\qfourInfo(p)}+\tfrac{1}{\qfourInfo(q)}$,
as required (interpreting $\tfrac1{\qfourInfo(p)}=+\infty$). By symmetry the case
$\qfourInfo(q)=0$ is identical.
\end{qfourobservation}

\section*{Remarks}

The proof mirrors Blachman's \cite{Blachman} derivation of the classical Stam
inequality \cite{Stam}: there the score of a sum is a conditional expectation of
the marginal score and the inequality reduces to Cauchy--Schwarz. Here the role
of that conditional-expectation step is played by
Observation~\ref{obs:scores} (the score of $p\qfourbn q$ is the image of the marginal
scores under the Jacobian $\qfourJac$) together with the operator-norm bound
Proposition~\ref{prop:contraction}. Real-rootedness of $\qfourbn$ is used exactly
once and essentially — through the hyperbolicity of \eqref{eq:walsh} and the
BGLS convexity theorem — which is why the inequality fails for convolutions that
do not preserve real-rootedness.

\begin{thebibliography}{9}
\bibitem{MSS22} A.~W.~Marcus, D.~A.~Spielman, N.~Srivastava.
  \emph{Finite free convolutions of polynomials}.
  Probab.\ Theory Related Fields \textbf{182} (2022), 807--848.
\bibitem{Walsh} J.~L.~Walsh.
  \emph{On the location of the roots of certain types of polynomials}.
  Trans.\ Amer.\ Math.\ Soc.\ \textbf{24} (1922), 163--180.
\bibitem{RS} Q.~I.~Rahman, G.~Schmeisser.
  \emph{Analytic Theory of Polynomials}.
  London Math.\ Soc.\ Monographs (N.S.) \textbf{26}, Oxford Univ.\ Press, 2002.
  (Grace--Walsh--Szeg\H{o} coincidence theorem: Thm.~3.4.1b.)
\bibitem{Nuij} W.~Nuij.
  \emph{A note on hyperbolic polynomials}.
  Math.\ Scand.\ \textbf{23} (1968), 69--72.
\bibitem{Blachman} N.~Blachman.
  \emph{The convolution inequality for entropy powers}.
  IEEE Trans.\ Inform.\ Theory \textbf{11} (1965), 267--271.
\bibitem{Garding} L.~G\aa rding.
  \emph{An inequality for hyperbolic polynomials}.
  J.\ Math.\ Mech.\ \textbf{8} (1959), 957--965.
\bibitem{BGLS} H.~H.~Bauschke, O.~G\"uler, A.~S.~Lewis, H.~S.~Sendov.
  \emph{Hyperbolic polynomials and convex analysis}.
  Canad.\ J.\ Math.\ \textbf{53} (2001), 470--488.
\bibitem{Whitney} H.~Whitney.
  \emph{Complex Analytic Varieties}.
  Addison-Wesley, Reading, MA, 1972. (For continuity of roots in the
  coefficients; equivalently a standard consequence of Rouch\'e's theorem.)
\bibitem{Stam} A.~Stam.
  \emph{Some inequalities satisfied by the quantities of information of Fisher
  and Shannon}. Inform.\ Control \textbf{2} (1959), 101--112.
\end{thebibliography}

\endgroup